% Complete reviewed Unit 6; editorial master: chapters/ch06/student.md
% Markdown SHA-256: 345c6edbcd113e28368865a2f456416bb2964a941ef097fb9bcc2906d037a545
% One-use initial assembly followed by complete content and page review.
\newcommand{\runningtitle}{Unit 6. Irregular verbs}
% Layout inherited from the reviewed previous course; set a running title per unit.
% This is a small hand-maintained TeX layout, not a conversion pipeline.
\RequirePackage{fix-cm}
\documentclass[11pt,twoside]{article}
\usepackage{fontspec}
\setmainfont{Cambria}[
  Renderer=HarfBuzz,
  BoldFont=cambriab.ttf,
  ItalicFont=cambriai.ttf,
  BoldItalicFont=cambriaz.ttf
]
\setmonofont{lmmono10-regular.otf}[BoldFont=lmmonolt10-bold.otf]
\usepackage[paperwidth=165mm,paperheight=241mm,inner=21mm,outer=19mm,
  top=22mm,bottom=20mm,headheight=14pt,headsep=7mm,footskip=10mm]{geometry}
\usepackage{xcolor,array,booktabs,longtable,tabularx}
\usepackage{enumitem,fancyhdr}
\usepackage[unicode,hidelinks]{hyperref}
\definecolor{ink}{HTML}{000000}
\definecolor{muted}{HTML}{4C5558}
\setlength{\parindent}{0pt}
\setlength{\parskip}{3.5pt}
\AtBeginDocument{\fontsize{11}{14.3}\selectfont}
\setlength{\tabcolsep}{4pt}
\renewcommand{\arraystretch}{1.10}
\setlist[itemize]{leftmargin=1.3em,label=\textbullet,itemsep=2pt,topsep=3pt,parsep=0pt}
\setlist[enumerate]{leftmargin=1.8em,itemsep=3pt,topsep=4pt,parsep=0pt}
\widowpenalty=10000
\clubpenalty=10000
\emergencystretch=1em
\raggedbottom
\pagestyle{fancy}
\fancyhf{}
\providecommand{\runningtitle}{Old English}
\fancyhead[LE,RO]{\small\itshape \runningtitle}
\fancyhead[LO,RE]{\small Old English}
\fancyfoot[LE,RO]{\small\thepage}
\renewcommand{\headrulewidth}{0.3pt}
\newcommand{\booktitle}[2]{
  {\fontsize{9}{11}\selectfont\color{muted}#1\par}
  \vspace{3pt}
  {\fontsize{20}{24}\selectfont\color{ink}#2\par}
  \vspace{6pt}
}
\newcommand{\source}[1]{\par{\footnotesize\color{muted}#1\par}}
\newcommand{\readingtext}[1]{
  \par\addvspace{5pt}
  \begingroup\leftskip=1.5em
    \fontsize{12}{16}\selectfont #1\par
  \endgroup\addvspace{5pt}
}
\makeatletter
\renewcommand{\section}{\@startsection{section}{1}{\z@}
  {-2.4ex \@plus -0.6ex \@minus -0.2ex}{0.8ex \@plus 0.2ex}
  {\normalfont\fontsize{12}{15}\selectfont\bfseries\color{ink}}}
\renewcommand{\subsection}{\@startsection{subsection}{2}{\z@}
  {-1.8ex \@plus -0.5ex \@minus -0.2ex}{0.5ex \@plus 0.2ex}
  {\normalfont\fontsize{11}{14}\selectfont\bfseries\color{ink}}}
\makeatother

\hyphenpenalty=10000
\exhyphenpenalty=10000
\emergencystretch=3em
\newcommand{\itempara}[2]{\par\addvspace{3pt}\noindent\hangindent=1.8em\hangafter=1\makebox[1.8em][l]{#1.}#2\par}
\newcommand{\reserve}[1]{\par\penalty-100\begingroup\dimen0=\pagegoal\advance\dimen0 by-\pagetotal\ifdim\dimen0<#1\newpage\fi\endgroup}
\begin{document}
\addfontfeatures{Ligatures=TeXOff}
\raggedright
\booktitle{UNIT 6}{Frequent irregular verbs, infinitives, and commands}

% mdparagraph:P001:begin
This unit adds frequent irregular indicative systems, commands addressed to one or more people, and infinitives used after modals or to express purpose.
% mdparagraph:P001:end

% mdparagraph:P002:begin
In the tables, \textbf{1}, \textbf{2}, and \textbf{3} mean first, second, and third person; \textbf{sg.} and \textbf{pl.} mean singular and plural. \textbf{1–3 pl.} gives the same form for all three plural persons. Dual subjects take plural finite verbs. Preserve vowel length and the distinction between e and ę. Examples and exercises are new teaching material unless a source is explicitly identified.
% mdparagraph:P002:end

\reserve{7\baselineskip}
% mdparagraph:P003:begin
\section*{1. Words and meanings}
% mdparagraph:P003:end

% mdparagraph:P004:begin
The verb labels in the first column identify systems. Some are conventional dictionary labels rather than usable infinitives; §2 explains the distinction. Learn the meanings given here with their stated constructions.
% mdparagraph:P004:end

% teaching: NEW-U06-VOC-witan
% teaching: NEW-U06-VOC-cunnan
% teaching: NEW-U06-VOC-thurfan
% teaching: NEW-U06-VOC-durran
% teaching: NEW-U06-VOC-sculan
% teaching: NEW-U06-VOC-magan
% teaching: NEW-U06-VOC-motan
% teaching: NEW-U06-VOC-willan
% teaching: NEW-U06-VOC-don
% teaching: NEW-U06-VOC-gan
% teaching: NEW-U06-VOC-thing-matters
% teaching: NEW-U06-VOC-to
% mdtable:T01:begin
\par\begingroup
\fontsize{10.5}{13.2}\selectfont
\setlength{\LTleft}{0pt}\setlength{\LTright}{\fill}
\noindent
\begin{tabular}{@{}>{\raggedright\arraybackslash}p{26.96mm}>{\raggedright\arraybackslash}p{34.04mm}>{\raggedright\arraybackslash}p{57.37mm}@{}}
\toprule
\textbf{Word or verb label} & \textbf{Meaning in this unit} & \textbf{Construction} \\
\midrule
witan & know & accusative fact or matter known \\
cunnan & know how; be able through skill & plain infinitive \\
ðurfan & need to & plain infinitive \\
durran & dare to & plain infinitive \\
sculan & must; be obliged to & plain infinitive \\
mugan, also called magan & can; be able to & plain infinitive \\
mōtan & may; be allowed to & plain infinitive \\
willan & wish; want to & plain infinitive \\
dōn & do; perform & accusative work or thing done \\
gān & go; walk & ordinary intransitive use \\
ðing, neuter & thing; a fact or matter known & the earlier noun forms, with this additional meaning \\
tō & to & before an inflected infinitive \\
\bottomrule
\end{tabular}
\par\endgroup

% mdtable:T01:end
% mdparagraph:P005:begin
The new sense of \textbf{ðing} includes a matter known or discussed. It does not replace the earlier sense of a thing or object. No new noun declension is needed. Infinitival tō is distinguished from locally supplied destination phrases in the source reading.
% mdparagraph:P005:end

\reserve{7\baselineskip}
% mdparagraph:P006:begin
\section*{2. Complete indicative systems}
% mdparagraph:P006:end

% teaching: NEW-U06-GRAM-PRETERITE-PRESENT-INDICATIVE
\reserve{7\baselineskip}
% mdparagraph:P007:begin
\subsection*{Preterite-present verbs}
% mdparagraph:P007:end

% mdparagraph:P008:begin
These verbs have present forms inherited from an older strong preterite. Thus \textbf{wāt} and \textbf{cann} resemble strong past forms but have present meaning. Their own preterites have a dental element, as weak preterites do. The historical name \textbf{preterite-present} does not make their present tense past in translation.
% mdparagraph:P008:end

% mdparagraph:P009:begin
The first and third singular present coincide. The second singular has its own form, and all three plural persons share another. Learn the singular/\allowbreak{}plural stem changes and the dental preterite together.
% mdparagraph:P009:end

\reserve{12\baselineskip}
% mdparagraph:P010:begin
\textbf{Present indicative.}
% mdparagraph:P010:end

% teaching: NEW-U06-PAR-witan
% teaching: NEW-U06-PAR-cunnan
% teaching: NEW-U06-PAR-thurfan
% teaching: NEW-U06-PAR-durran
% teaching: NEW-U06-PAR-sculan
% teaching: NEW-U06-PAR-magan
% teaching: NEW-U06-PAR-motan
% table: NEW-U06-DISPLAY-PRETERITE-PRESENTS-PRES
% mdtable:T02:begin
% course-table-style: 2026-10
% course-panel:T02-P01:begin
\par\begingroup
\fontsize{10.5}{13.2}\selectfont
\setlength{\LTleft}{0pt}\setlength{\LTright}{\fill}
\noindent
\begin{tabular}{@{}>{\raggedright\arraybackslash}p{24.232mm}>{\raggedright\arraybackslash}p{22.13mm}>{\raggedright\arraybackslash}p{22.13mm}>{\raggedright\arraybackslash}p{22.13mm}>{\raggedright\arraybackslash}p{22.13mm}@{}}
\toprule
\textbf{Person and number} & \textbf{witan} & \textbf{cunnan} & \textbf{ðurfan} & \textbf{durran} \\
\midrule
1 sg. & wāt & cann & ðearf & dearr \\
2 sg. & wāst & canst & ðearft & dearst \\
3 sg. & wāt & cann & ðearf & dearr \\
1–3 pl. & witon & cunnon & ðurfon & durron \\
\bottomrule
\end{tabular}
\par\endgroup
% course-panel:T02-P01:end
\par\addvspace{5pt}
% course-panel:T02-P02:begin
\par\begingroup
\fontsize{10.5}{13.2}\selectfont
\setlength{\LTleft}{0pt}\setlength{\LTright}{\fill}
\noindent
\begin{tabular}{@{}>{\raggedright\arraybackslash}p{30.902mm}>{\raggedright\arraybackslash}p{28.221mm}>{\raggedright\arraybackslash}p{28.221mm}>{\raggedright\arraybackslash}p{28.221mm}@{}}
\toprule
\textbf{Person and number} & \textbf{sculan} & \textbf{mugan} & \textbf{mōtan} \\
\midrule
1 sg. & sceal & mæg & mōt \\
2 sg. & scealt & meaht & mōst \\
3 sg. & sceal & mæg & mōt \\
1–3 pl. & sculon & magon & mōton \\
\bottomrule
\end{tabular}
\par\endgroup
% course-panel:T02-P02:end
% mdtable:T02:end
\reserve{12\baselineskip}
% mdparagraph:P011:begin
\textbf{Preterite indicative.}
% mdparagraph:P011:end

% table: NEW-U06-DISPLAY-PRETERITE-PRESENTS-PRET
% mdtable:T03:begin
% course-table-style: 2026-10
% course-panel:T03-P01:begin
\par\begingroup
\fontsize{10.5}{13.2}\selectfont
\setlength{\LTleft}{0pt}\setlength{\LTright}{\fill}
\noindent
\begin{tabular}{@{}>{\raggedright\arraybackslash}p{24.232mm}>{\raggedright\arraybackslash}p{22.13mm}>{\raggedright\arraybackslash}p{22.13mm}>{\raggedright\arraybackslash}p{22.13mm}>{\raggedright\arraybackslash}p{22.13mm}@{}}
\toprule
\textbf{Person and number} & \textbf{witan} & \textbf{cunnan} & \textbf{ðurfan} & \textbf{durran} \\
\midrule
1 sg. & wiste & cūðe & ðorfte & dorste \\
2 sg. & wistest & cūðest & ðorftest & dorstest \\
3 sg. & wiste & cūðe & ðorfte & dorste \\
1–3 pl. & wiston & cūðon & ðorfton & dorston \\
\bottomrule
\end{tabular}
\par\endgroup
% course-panel:T03-P01:end
\par\addvspace{5pt}
% course-panel:T03-P02:begin
\par\begingroup
\fontsize{10.5}{13.2}\selectfont
\setlength{\LTleft}{0pt}\setlength{\LTright}{\fill}
\noindent
\begin{tabular}{@{}>{\raggedright\arraybackslash}p{30.902mm}>{\raggedright\arraybackslash}p{28.221mm}>{\raggedright\arraybackslash}p{28.221mm}>{\raggedright\arraybackslash}p{28.221mm}@{}}
\toprule
\textbf{Person and number} & \textbf{sculan} & \textbf{mugan} & \textbf{mōtan} \\
\midrule
1 sg. & sceolde & meahte & mōste \\
2 sg. & sceoldest & meahtest & mōstest \\
3 sg. & sceolde & meahte & mōste \\
1–3 pl. & sceoldon & meahton & mōston \\
\bottomrule
\end{tabular}
\par\endgroup
% course-panel:T03-P02:end
% mdtable:T03:end
% mdparagraph:P012:begin
For each verb, the present and preterite tables give all twelve person–number–tense cells. In the preterite, the first and third singular end in \textbf{-e}; the second singular adds \textbf{-st}, and the plural changes \textbf{-e} to \textbf{-on}. Keep each verb’s printed stem. The second singular here does not follow the strong-preterite pattern learned in Units 4–5.
% mdparagraph:P012:end

% mdparagraph:P013:begin
\textbf{Selected forms and alternatives.} This course selects witan, witon, and wiste. Wieton/\allowbreak{}wioton and the wisse-series are documented alternatives. Cook also gives can beside cann and scolde beside sceolde. For mugan, the tables select meaht and meahte; Wright records miht and mehte/\allowbreak{}mihte. Use the selected forms when an exercise asks for them; the alternatives are not errors merely because they differ.
% mdparagraph:P013:end

% teaching: NEW-U06-GRAM-DEFECTIVE-REFERENCE
% mdparagraph:P014:begin
\textbf{Citation labels and missing forms.} Cook’s \textbf{mugan} and the common label \textbf{magan} name the same system. Fulk marks the conventional infinitives *durran, *sculan, *mōtan, and *magan as unattested. Cook marks several labels as doubtful. A label ending in -an is therefore not sufficient evidence for an infinitive that can be put into a sentence. Do not make an inflected infinitive or a command by adding ordinary endings to these labels.
% mdparagraph:P014:end

% mdparagraph:P015:begin
The labels witan, cunnan, and ðurfan are attested uninflected infinitives, but their infinitival uses and complete command systems are not taught here. Learn the finite systems above and the selected infinitive pairs below. A missing cell is different from two genuine cells that happen to share a spelling, such as the first and third singular wāt.
% mdparagraph:P015:end

% mdparagraph:P016:begin
Cook and Wright call certain witan forms commands; Fulk describes such preterite-present command uses as subjunctive substitutes. They do not establish an ordinary imperative pair for every modal. These substitutes and first- or third-person wishes are outside this unit’s command practice.
% mdparagraph:P016:end

% mdparagraph:P017:begin
\textbf{Sources.} Albert S. Cook, \textbf{A First Book in Old English: Grammar, Reader, Notes, and Vocabulary}, second edition (1895), §§125–137, pp.76–78 /\allowbreak{} PDF pp.92–94. R. D. Fulk, \textbf{An Introductory Grammar of Old English with an Anthology of Readings}, corrected 2020 reprint, §123, pp.89–90 /\allowbreak{} PDF pp.101–102, distinguishes unattested labels and command substitutes. Joseph and Elizabeth Mary Wright, \textbf{Old English Grammar}, second edition (1914), gives the corroborating witan and magan variants; see the precise variant references in §9.
% mdparagraph:P017:end

\reserve{7\baselineskip}
% mdparagraph:P018:begin
\subsection*{Dōn, gān, and willan}
% mdparagraph:P018:end

% teaching: NEW-U06-GRAM-DON-INDICATIVE
% teaching: NEW-U06-GRAM-GAN-INDICATIVE
% teaching: NEW-U06-GRAM-WILLAN-INDICATIVE
% mdparagraph:P019:begin
Learn each present with its preterite. \textbf{Dōn} changes ō to ē in dēst and dēð, and has the past dyde. \textbf{Gān} changes ā to ǣ in gǣst and gǣð; its past \textbf{ēode} comes from a different stem. \textbf{Willan} has wille in both the first and third singular present, with wilt in the second singular. Do not add an ordinary third-singular -ð to wille.
% mdparagraph:P019:end

\reserve{8\baselineskip}
% mdparagraph:P020:begin
\textbf{Present indicative.}
% mdparagraph:P020:end

% teaching: NEW-U06-PAR-willan
% teaching: NEW-U06-PAR-don
% teaching: NEW-U06-PAR-gan
% table: NEW-U06-DISPLAY-ANOMALOUS-PRES
% mdtable:T04:begin
% course-table-style: 2026-10
% course-panel:T04-P01:begin
\par\begingroup
\fontsize{10.5}{13.2}\selectfont
\setlength{\LTleft}{0pt}\setlength{\LTright}{\fill}
\noindent
\begin{tabular}{@{}>{\raggedright\arraybackslash}p{30.902mm}>{\raggedright\arraybackslash}p{28.221mm}>{\raggedright\arraybackslash}p{28.221mm}>{\raggedright\arraybackslash}p{28.221mm}@{}}
\toprule
\textbf{Person and number} & \textbf{dōn} & \textbf{gān} & \textbf{willan} \\
\midrule
1 sg. & dō & gā & wille \\
2 sg. & dēst & gǣst & wilt \\
3 sg. & dēð & gǣð & wille \\
1–3 pl. & dōð & gāð & willað \\
\bottomrule
\end{tabular}
\par\endgroup
% course-panel:T04-P01:end
% mdtable:T04:end
\reserve{8\baselineskip}
% mdparagraph:P021:begin
\textbf{Preterite indicative.}
% mdparagraph:P021:end

% table: NEW-U06-DISPLAY-ANOMALOUS-PRET
% mdtable:T05:begin
% course-table-style: 2026-10
% course-panel:T05-P01:begin
\par\begingroup
\fontsize{10.5}{13.2}\selectfont
\setlength{\LTleft}{0pt}\setlength{\LTright}{\fill}
\noindent
\begin{tabular}{@{}>{\raggedright\arraybackslash}p{30.902mm}>{\raggedright\arraybackslash}p{28.221mm}>{\raggedright\arraybackslash}p{28.221mm}>{\raggedright\arraybackslash}p{28.221mm}@{}}
\toprule
\textbf{Person and number} & \textbf{dōn} & \textbf{gān} & \textbf{willan} \\
\midrule
1 sg. & dyde & ēode & wolde \\
2 sg. & dydest & ēodest & woldest \\
3 sg. & dyde & ēode & wolde \\
1–3 pl. & dydon & ēodon & woldon \\
\bottomrule
\end{tabular}
\par\endgroup
% course-panel:T05-P01:end
% mdtable:T05:end
% mdparagraph:P022:begin
Cook also gives \textbf{wile} beside wille. The table selects the spelling with double l. The negative contracted companions are not additional paradigms to learn in this unit.
% mdparagraph:P022:end

% mdparagraph:P023:begin
\textbf{Sources.} Cook §§138–141, pp.78–80 /\allowbreak{} PDF pp.94–96. Fulk §134, p.102 /\allowbreak{} PDF p.114, corroborates the complete dōn, gān, and willan indicative systems.
% mdparagraph:P023:end

\reserve{7\baselineskip}
% mdparagraph:P024:begin
\subsection*{The alternative bēo present}
% mdparagraph:P024:end

% teaching: NEW-U06-GRAM-BEON-INDICATIVE
% mdparagraph:P025:begin
The following alternative present series belongs with the earlier verb “be.” Both present series share the familiar preterite: wæs in the first and third singular, wǣre in the second singular, and wǣron in all plural persons. There is no separate bēon past to construct.
% mdparagraph:P025:end

% teaching: NEW-U06-PAR-wesan-beon
% table: NEW-U06-DISPLAY-BEON
% mdtable:T06:begin
% course-table-style: 2026-10
% course-panel:T06-P01:begin
\par\begingroup
\fontsize{10.5}{13.2}\selectfont
\setlength{\LTleft}{0pt}\setlength{\LTright}{\fill}
\noindent
\begin{tabular}{@{}>{\raggedright\arraybackslash}p{39.459mm}>{\raggedright\arraybackslash}p{39.459mm}>{\raggedright\arraybackslash}p{39.459mm}@{}}
\toprule
\textbf{Person and number} & \textbf{Earlier eom series} & \textbf{Alternative bēo series} \\
\midrule
1 sg. & eom & bēo \\
2 sg. & eart & bist \\
3 sg. & is & bið \\
1–3 pl. & sindon & bēoð \\
\bottomrule
\end{tabular}
\par\endgroup
% course-panel:T06-P01:end
% mdtable:T06:end
% teaching: NEW-U06-GRAM-BEON-USE
% mdparagraph:P026:begin
The bēo series can describe an expected later state. The predicate still follows the agreement already learned. In the following example the ship’s readiness is expected later.
% mdparagraph:P026:end

% example: NEW-U06-EXAMPLE-09
% mdparagraph:P027:begin
Ðæt scip bið gearu.
% mdparagraph:P027:end

% mdparagraph:P028:begin
The ship will be ready, in an explicitly expected later situation.
% mdparagraph:P028:end

% mdparagraph:P029:begin
Bēo forms also occur with habitual or general truths. Context matters: the two present series are not an absolute division into “now” and “future,” and every bið need not be translated “will be.” The exercises here give an explicit expected-later context.
% mdparagraph:P029:end

% mdparagraph:P030:begin
\textbf{Sources.} Cook §138, pp.78–79 /\allowbreak{} PDF pp.94–95; Fulk §52, p.33 /\allowbreak{} PDF p.45.
% mdparagraph:P030:end

\reserve{7\baselineskip}
% mdparagraph:P031:begin
\subsection*{Practice A. Indicative systems and citation labels}
% mdparagraph:P031:end

\reserve{5\baselineskip}
\par\begin{minipage}{\linewidth}\raggedright
% mdparagraph:P032:begin
\textbf{A1.}
% mdparagraph:P032:end

% item: NEW-U06-A01
% mdparagraph:P033:begin
Give all six person-and-number cells of the present indicative and all six of the preterite indicative of witan. Use the selected course forms and label every cell even where spellings coincide.
% mdparagraph:P033:end

\end{minipage}\par\addvspace{4pt}

\par\begin{minipage}{\linewidth}\raggedright
% mdparagraph:P034:begin
\textbf{A2.}
% mdparagraph:P034:end

% item: NEW-U06-A02
% mdparagraph:P035:begin
Give all six person-and-number cells of the present indicative and all six of the preterite indicative of cunnan. Use the selected course forms and label every cell even where spellings coincide.
% mdparagraph:P035:end

\end{minipage}\par\addvspace{4pt}

\par\begin{minipage}{\linewidth}\raggedright
% mdparagraph:P036:begin
\textbf{A3.}
% mdparagraph:P036:end

% item: NEW-U06-A03
% mdparagraph:P037:begin
Give all six person-and-number cells of the present indicative and all six of the preterite indicative of ðurfan. Use the selected course forms and label every cell even where spellings coincide.
% mdparagraph:P037:end

\end{minipage}\par\addvspace{4pt}

\par\begin{minipage}{\linewidth}\raggedright
% mdparagraph:P038:begin
\textbf{A4.}
% mdparagraph:P038:end

% item: NEW-U06-A04
% mdparagraph:P039:begin
Give all six person-and-number cells of the present indicative and all six of the preterite indicative of durran. Use the selected course forms and label every cell even where spellings coincide.
% mdparagraph:P039:end

\end{minipage}\par\addvspace{4pt}

\par\begin{minipage}{\linewidth}\raggedright
% mdparagraph:P040:begin
\textbf{A5.}
% mdparagraph:P040:end

% item: NEW-U06-A05
% mdparagraph:P041:begin
Give all six person-and-number cells of the present indicative and all six of the preterite indicative of sculan. Use the selected course forms and label every cell even where spellings coincide.
% mdparagraph:P041:end

\end{minipage}\par\addvspace{4pt}

\par\begin{minipage}{\linewidth}\raggedright
% mdparagraph:P042:begin
\textbf{A6.}
% mdparagraph:P042:end

% item: NEW-U06-A06
% mdparagraph:P043:begin
Give all six person-and-number cells of the present indicative and all six of the preterite indicative of mugan. Use the selected course forms and label every cell even where spellings coincide.
% mdparagraph:P043:end

\end{minipage}\par\addvspace{4pt}

\par\begin{minipage}{\linewidth}\raggedright
% mdparagraph:P044:begin
\textbf{A7.}
% mdparagraph:P044:end

% item: NEW-U06-A07
% mdparagraph:P045:begin
Give all six person-and-number cells of the present indicative and all six of the preterite indicative of mōtan. Use the selected course forms and label every cell even where spellings coincide.
% mdparagraph:P045:end

\end{minipage}\par\addvspace{4pt}

\par\begin{minipage}{\linewidth}\raggedright
% mdparagraph:P046:begin
\textbf{A8.}
% mdparagraph:P046:end

% item: NEW-U06-A08
% mdparagraph:P047:begin
Give all six person-and-number cells of the present indicative and all six of the preterite indicative of dōn. Use the selected course forms and label every cell even where spellings coincide.
% mdparagraph:P047:end

\end{minipage}\par\addvspace{4pt}

\par\begin{minipage}{\linewidth}\raggedright
% mdparagraph:P048:begin
\textbf{A9.}
% mdparagraph:P048:end

% item: NEW-U06-A09
% mdparagraph:P049:begin
Give all six person-and-number cells of the present indicative and all six of the preterite indicative of gān. Use the selected course forms and label every cell even where spellings coincide.
% mdparagraph:P049:end

\end{minipage}\par\addvspace{4pt}

\par\begin{minipage}{\linewidth}\raggedright
% mdparagraph:P050:begin
\textbf{A10.}
% mdparagraph:P050:end

% item: NEW-U06-A10
% mdparagraph:P051:begin
Give all six person-and-number cells of the present indicative and all six of the preterite indicative of willan. Use the selected course forms and label every cell even where spellings coincide.
% mdparagraph:P051:end

\end{minipage}\par\addvspace{4pt}

\par\begin{minipage}{\linewidth}\raggedright
% mdparagraph:P052:begin
\textbf{A11.}
% mdparagraph:P052:end

% item: NEW-U06-A11
% mdparagraph:P053:begin
Give all six person-and-number cells of the bēo present indicative series. Then give the corresponding earlier eom series. Does bēo supply a separate past-tense paradigm?
% mdparagraph:P053:end

\end{minipage}\par\addvspace{4pt}

\par\begin{minipage}{\linewidth}\raggedright
% mdparagraph:P054:begin
\textbf{A12.}
% mdparagraph:P054:end

% item: NEW-U06-A12
% mdparagraph:P055:begin
Compare wāt and wiste; cann and cūðe; gǣð and ēode; dēð and dyde. Using a third-person singular subject, identify the tense of each. Which forms can also serve a first-person singular subject? Explain why wāt and cann are not translated as past-tense forms.
% mdparagraph:P055:end

\end{minipage}\par\addvspace{4pt}

\par\begin{minipage}{\linewidth}\raggedright
% mdparagraph:P056:begin
\textbf{A13.}
% mdparagraph:P056:end

% item: NEW-U06-A13
% mdparagraph:P057:begin
The reference notes mark some modal infinitive citations as unattested or uncertain. Explain why a dictionary label such as mugan does not authorize you to create an infinitive or a command for every modal verb. Distinguish a missing paradigm cell from two attested cells with the same spelling.
% mdparagraph:P057:end

\end{minipage}\par\addvspace{4pt}

\reserve{7\baselineskip}
% mdparagraph:P058:begin
\section*{3. Infinitives and their uses}
% mdparagraph:P058:end

% teaching: NEW-U06-GRAM-INFINITIVE-FORMS
% teaching: NEW-U06-INF-MARKER
% mdparagraph:P059:begin
An \textbf{infinitive} names an action without marking person or number. The \textbf{plain infinitive} normally ends in -an or -ian, with contracted forms such as gān and dōn. In earlier units many infinitives were printed only to name verbs. The forms in this unit’s tables are now learned for the stated uses.
% mdparagraph:P059:end

% mdparagraph:P060:begin
The \textbf{inflected infinitive} follows the separate word \textbf{tō}. For the selected -an forms, replace -an with \textbf{-anne}; -ian consequently gives \textbf{-ianne}. Keep the stem consonants: doubled consonants remain doubled in the inflected infinitive. Contracted forms ending in -n add -ne, giving \textbf{-nne}: the complete forms appear in §4. The infinitive is not conjugated for the subject’s person or number.
% mdparagraph:P060:end

% mdparagraph:P061:begin
The tables select \textbf{-anne}. \textbf{-enne} is also a documented ending; it is a valid variant, although an exercise expressly requesting the selected -anne form still asks for that form. Infinitive formation does not supply missing modal forms or teach participles, perfects, or passives.
% mdparagraph:P061:end

\reserve{7\baselineskip}
% mdparagraph:P062:begin
\subsection*{A finite modal with a plain infinitive}
% mdparagraph:P062:end

% teaching: NEW-U06-GRAM-PLAIN-INFINITIVE-COMPLEMENT
% mdparagraph:P063:begin
In the constructions below, the finite verb agrees with the subject and marks present or past. The plain infinitive expresses the action performed by that same subject. It takes neither personal endings nor tō in this construction. Two or more people, including a dual-pronoun subject, take a plural finite verb; the infinitive itself stays unchanged.
% mdparagraph:P063:end

% teaching: NEW-U06-GRAM-MODAL-CUNNAN
% teaching: NEW-U06-GRAM-MODAL-MAGAN
% mdparagraph:P064:begin
\textbf{Cunnan} expresses knowing how through learning or skill. \textbf{Mugan} expresses capability: the subject is able to perform the action. English “can” can represent either, so select the intended sense from the context.
% mdparagraph:P064:end

% example: NEW-U06-EXAMPLE-01
% mdparagraph:P065:begin
Ic cann swimman. Hē mæg ðæt wæter beran.
% mdparagraph:P065:end

% mdparagraph:P066:begin
I know how to swim. He can carry the water.
% mdparagraph:P066:end

% mdparagraph:P067:begin
For the skill statements, learn these additional infinitive pairs. No imperative pair for either verb is required here.
% mdparagraph:P067:end

% teaching: NEW-U06-NONFINITE-singan
% teaching: NEW-U06-NONFINITE-swimman
% table: NEW-U06-DISPLAY-ADDITIONAL-INFINITIVES
% mdtable:T07:begin
\par\begingroup
\fontsize{10.5}{13.2}\selectfont
\setlength{\LTleft}{0pt}\setlength{\LTright}{\fill}
\noindent
\begin{tabular}{@{}>{\raggedright\arraybackslash}p{28.7mm}>{\raggedright\arraybackslash}p{36.45mm}>{\raggedright\arraybackslash}p{53.23mm}@{}}
\toprule
\textbf{Verb} & \textbf{Plain infinitive} & \textbf{Inflected infinitive after tō} \\
\midrule
singan & singan & singanne \\
swimman & swimman & swimmanne \\
\bottomrule
\end{tabular}
\par\endgroup

% mdtable:T07:end
% teaching: NEW-U06-GRAM-MODAL-THURFAN
% teaching: NEW-U06-GRAM-MODAL-DURRAN
% teaching: NEW-U06-GRAM-MODAL-SCULAN
% teaching: NEW-U06-GRAM-MODAL-MOTAN
% mdparagraph:P068:begin
\textbf{Ðurfan} expresses need to do something; \textbf{durran}, daring to do it. \textbf{Sculan} expresses obligation or necessity in the selected use. \textbf{Mōtan} expresses permission, “be allowed to.” A past form can describe actual past need, daring, obligation, or permission: the morphology alone does not make the statement hypothetical.
% mdparagraph:P068:end

% mdparagraph:P069:begin
In particular, selected \textbf{mōste} can mean “was allowed to”; it need not mean modern English “must.” Other senses occur in the language, but the exercises state the sense required here.
% mdparagraph:P069:end

% example: NEW-U06-EXAMPLE-02
% mdparagraph:P070:begin
Wē sculon ðæt weorc gefręmman.
% mdparagraph:P070:end

% mdparagraph:P071:begin
We must perform the work.
% mdparagraph:P071:end

% teaching: NEW-U06-GRAM-VOLITION-WILLAN
% teaching: NEW-U06-GRAM-FARAN-INTRANSITIVE
% mdparagraph:P072:begin
\textbf{Willan} expresses wishing or wanting. Its preterite \textbf{wolde} can state past desire. \textbf{Faran} is used intransitively for going or travelling; it can name the desired action without an expressed object or destination.
% mdparagraph:P072:end

% example: NEW-U06-EXAMPLE-03
% mdparagraph:P073:begin
Hē wolde faran.
% mdparagraph:P073:end

% mdparagraph:P074:begin
He wanted to travel.
% mdparagraph:P074:end

% mdparagraph:P075:begin
Wanting to travel does not by itself assert that the journey occurred. Do not translate every wolde as hypothetical “would,” or treat willan as an automatic future-tense auxiliary.
% mdparagraph:P075:end

% teaching: NEW-U06-GRAM-MODAL-NEGATION
% mdparagraph:P076:begin
The familiar negative \textbf{ne} precedes the finite verb. Use this ordinary pattern with the selected durran and mugan clauses. Contracted negative paradigms and negative concord are not needed.
% mdparagraph:P076:end

\reserve{7\baselineskip}
% mdparagraph:P077:begin
\subsection*{Witan with a noun object}
% mdparagraph:P077:end

% teaching: NEW-U06-GRAM-WITAN-OBJECT
% mdparagraph:P078:begin
\textbf{Witan}, “know,” takes an accusative nominal object for the matter known. The following “things” are facts or matters already discussed. The noun’s form is familiar; its fact/\allowbreak{}matter meaning is new here.
% mdparagraph:P078:end

% example: NEW-U06-EXAMPLE-10
% mdparagraph:P079:begin
Ic wāt ðā ðing.
% mdparagraph:P079:end

% mdparagraph:P080:begin
I know those things, referring to matters already discussed.
% mdparagraph:P080:end

% mdparagraph:P081:begin
This is knowledge of something, whereas cunnan with an infinitive describes knowing how to do something. The object here is a noun phrase; no finite “that” clause is required.
% mdparagraph:P081:end

\reserve{7\baselineskip}
% mdparagraph:P082:begin
\subsection*{Ordinary action and purpose}
% mdparagraph:P082:end

% teaching: NEW-U06-GRAM-DON-OBJECT
% teaching: NEW-U06-GRAM-GAN-INTRANSITIVE
% teaching: NEW-U06-GRAM-CUMAN-INTRANSITIVE
% mdparagraph:P083:begin
Use \textbf{dōn}, “do; perform,” with an accusative work or thing done. Use \textbf{gān}, “go,” and \textbf{cuman}, “come,” in ordinary intransitive clauses. These uses do not require a destination phrase and do not establish a general rule for motion prepositions.
% mdparagraph:P083:end

% teaching: NEW-U06-GRAM-PURPOSE-INFINITIVE
% mdparagraph:P084:begin
After \textbf{cuman}, \textbf{tō + inflected infinitive} can state the purpose of coming. The person or people who come are also the understood performers of the infinitive’s action.
% mdparagraph:P084:end

% example: NEW-U06-EXAMPLE-07
% mdparagraph:P085:begin
Hē cōm tō helpanne.
% mdparagraph:P085:end

% mdparagraph:P086:begin
He came to help.
% mdparagraph:P086:end

% mdparagraph:P087:begin
Here he came, and he intended to help. The infinitive names the purpose; it does not assert that the helping was completed. This is the purpose construction practised below. Old English also permits certain bare infinitives after motion verbs, so it would be wrong to say that every grammatical purpose must have tō. When an exercise requests the inflected construction, provide that target.
% mdparagraph:P087:end

% teaching: NEW-U06-GRAM-INFINITIVE-GOVERNMENT
% mdparagraph:P088:begin
An infinitive keeps the object relations of its verb. Carrying, taking, and eating can have accusative objects; \textbf{helpan} takes the person helped in the dative. Giving has an accusative thing and a dative recipient. Adding a modal or changing to an inflected infinitive does not turn those datives into accusatives.
% mdparagraph:P088:end

% mdparagraph:P089:begin
\textbf{Sources.} Cook §§199–200, pp.105–106 /\allowbreak{} PDF pp.121–122; Fulk §58, p.38 /\allowbreak{} PDF p.50. Cook §200 provides cuman with an inflected purpose infinitive; §199.1 also records bare infinitives with motion and purpose. The modal distinctions follow \textbf{Bosworth–Toller, An Anglo-Saxon Dictionary}, main dictionary (1898): CUNNAN II, p.174b; MAGAN III(1), p.665a; ÞURFAN II(1), p.1076b; DURRAN, p.219b; SCULAN II(1), pp.843b–844; MÓTAN I, p.699a–b; WILLAN II(a), p.1226b. Ordinary DÓN is at p.208b, GÁN at p.360b, and FARAN I at p.270b. Witan’s nominal object is evidenced in Cook Reader XII, p.173 /\allowbreak{} PDF p.189; the known-fact sense of þing appears under WITAN I(6)(a), Bosworth–Toller p.1244b. Cuman is glossed in Cook’s vocabulary, p.282 /\allowbreak{} PDF p.298.
% mdparagraph:P089:end

\reserve{7\baselineskip}
% mdparagraph:P090:begin
\section*{4. Commands and complete model tables}
% mdparagraph:P090:end

% teaching: NEW-U06-GRAM-IMPERATIVE-USE
% mdparagraph:P091:begin
The \textbf{imperative} directly addresses the person who is to act. Use the second singular for one addressee and the second plural for two or more. There is no separate dual imperative ending. The subject may be omitted; an expressed subject can precede the verb. The commands used here omit it unless the exercise states otherwise.
% mdparagraph:P091:end

% example: NEW-U06-EXAMPLE-04
% mdparagraph:P092:begin
Nim ðone męte. Nimað ðone męte.
% mdparagraph:P092:end

% mdparagraph:P093:begin
Take the food, one addressee. Take the food, several addressees.
% mdparagraph:P093:end

% mdparagraph:P094:begin
The number of the command identifies the addressees, not the number of its object. A command concerning the same food may address one person or several.
% mdparagraph:P094:end

% mdparagraph:P095:begin
The following tables give each complete command pair beside its complete plain/\allowbreak{}inflected infinitive pair. \textbf{The last column is used after separate tō.} Learn every form printed here. Most plural commands match the present indicative plural; that shared spelling does not make their uses identical.
% mdparagraph:P095:end

\reserve{7\baselineskip}
% mdparagraph:P096:begin
\subsection*{Weak I}
% mdparagraph:P096:end

% teaching: NEW-U06-GRAM-WEAK-IMPERATIVE-FORMS
% mdparagraph:P097:begin
The heavy-stem model has no final -e in the singular command: \textbf{hīer}. The short and geminate patterns have \textbf{-e}, with the singular consonant pattern shown in \textbf{sęle, sęte, gefręme, nęre}. In nęre the i of nęrian is absent; it remains in the plural nęriað. \textbf{Gierwan} has \textbf{giere}, without w, beside gierwað.
% mdparagraph:P097:end

% mdparagraph:P098:begin
Learn \textbf{cys} and \textbf{gefyl} with a single final consonant, beside cyssað and gefyllað. Their selected final simplification does not turn them into the -e type. Conversely, sęttan has sęte; its double consonant is not a reason to use the cys pattern. Final consonant simplification is frequent, not a claim that all unsimplified final spellings are impossible.
% mdparagraph:P098:end

% teaching: NEW-U06-IMP-sellan
% teaching: NEW-U06-NONFINITE-sellan
% teaching: NEW-U06-IMP-hieran
% teaching: NEW-U06-NONFINITE-hieran
% teaching: NEW-U06-IMP-gefyllan
% teaching: NEW-U06-NONFINITE-gefyllan
% teaching: NEW-U06-IMP-cyssan
% teaching: NEW-U06-NONFINITE-cyssan
% teaching: NEW-U06-IMP-settan
% teaching: NEW-U06-NONFINITE-settan
% teaching: NEW-U06-IMP-gierwan
% teaching: NEW-U06-NONFINITE-gierwan
% teaching: NEW-U06-IMP-gefremman
% teaching: NEW-U06-NONFINITE-gefremman
% teaching: NEW-U06-IMP-nerian
% teaching: NEW-U06-NONFINITE-nerian
% table: NEW-U06-DISPLAY-WEAK-I
% mdtable:T08:begin
\par\begingroup
\fontsize{10.5}{13.2}\selectfont
\setlength{\LTleft}{0pt}\setlength{\LTright}{\fill}
\noindent
\begin{tabular}{@{}>{\raggedright\arraybackslash}p{18.5mm}>{\raggedright\arraybackslash}p{17.13mm}>{\raggedright\arraybackslash}p{25.07mm}>{\raggedright\arraybackslash}p{21.44mm}>{\raggedright\arraybackslash}p{30.62mm}@{}}
\toprule
\textbf{Verb} & \textbf{Command to one} & \textbf{Command to two or more} & \textbf{Plain infinitive} & \textbf{Inflected infinitive after tō} \\
\midrule
sęllan & sęle & sęllað & sęllan & sęllanne \\
hīeran & hīer & hīerað & hīeran & hīeranne \\
gefyllan & gefyl & gefyllað & gefyllan & gefyllanne \\
cyssan & cys & cyssað & cyssan & cyssanne \\
sęttan & sęte & sęttað & sęttan & sęttanne \\
gierwan & giere & gierwað & gierwan & gierwanne \\
gefręmman & gefręme & gefręmmað & gefręmman & gefręmmanne \\
nęrian & nęre & nęriað & nęrian & nęrianne \\
\bottomrule
\end{tabular}
\par\endgroup

% mdtable:T08:end
\reserve{7\baselineskip}
% mdparagraph:P099:begin
\subsection*{Weak II and III}
% mdparagraph:P099:end

% mdparagraph:P100:begin
Weak II changes \textbf{-ian} to \textbf{-a} for the singular command and retains \textbf{-iað} in the plural: lufa/\allowbreak{}lufiað. Its inflected infinitive has -ianne.
% mdparagraph:P100:end

% mdparagraph:P101:begin
The four learned weak-III verbs require their own singular forms: \textbf{hafa, leofa, saga, hoga}. Their plural commands use the familiar present-plural stems. Form habbanne from habban; the command stem hafa is different.
% mdparagraph:P101:end

% teaching: NEW-U06-IMP-lufian
% teaching: NEW-U06-NONFINITE-lufian
% teaching: NEW-U06-IMP-habban
% teaching: NEW-U06-NONFINITE-habban
% teaching: NEW-U06-IMP-libban
% teaching: NEW-U06-NONFINITE-libban
% teaching: NEW-U06-IMP-secgan
% teaching: NEW-U06-NONFINITE-secgan
% teaching: NEW-U06-IMP-hycgan
% teaching: NEW-U06-NONFINITE-hycgan
% table: NEW-U06-DISPLAY-WEAK-II-III
% mdtable:T09:begin
\par\begingroup
\fontsize{10.5}{13.2}\selectfont
\setlength{\LTleft}{0pt}\setlength{\LTright}{\fill}
\noindent
\begin{tabular}{@{}>{\raggedright\arraybackslash}p{17.31mm}>{\raggedright\arraybackslash}p{17.57mm}>{\raggedright\arraybackslash}p{23.8mm}>{\raggedright\arraybackslash}p{21.98mm}>{\raggedright\arraybackslash}p{32.1mm}@{}}
\toprule
\textbf{Verb} & \textbf{Command to one} & \textbf{Command to two or more} & \textbf{Plain infinitive} & \textbf{Inflected infinitive after tō} \\
\midrule
lufian & lufa & lufiað & lufian & lufianne \\
habban & hafa & habbað & habban & habbanne \\
libban & leofa & libbað & libban & libbanne \\
sęcgan & saga & sęcgað & sęcgan & sęcganne \\
hycgan & hoga & hycgað & hycgan & hycganne \\
\bottomrule
\end{tabular}
\par\endgroup

% mdtable:T09:end
% mdparagraph:P102:begin
Cook also gives \textbf{sęge} beside saga and \textbf{lifiað} beside libbað, and optional e in \textbf{sęcg(e)að} and \textbf{hycg(e)að}. The tables select saga, libbað, sęcgað, and hycgað. The earlier recognition-only negative companion nabban receives no new command requirement.
% mdparagraph:P102:end

% mdparagraph:P103:begin
\textbf{Sources for the weak forms.} Cook §§112–124, pp.68–76 /\allowbreak{} PDF pp.84–92. Final-geminate simplification is described in §35, pp.24–25 /\allowbreak{} PDF pp.40–41; cys and gefyl apply that rule to the heavy weak-I pattern, rather than being separate continuous-text quotations. Wright §§521–525 and 529–530, pp.266–268 and 270–271 /\allowbreak{} PDF pp.286–288 and 290–291, corroborates the short, geminate, and gierwan patterns.
% mdparagraph:P103:end

\reserve{7\baselineskip}
% mdparagraph:P104:begin
\subsection*{Strong I–III}
% mdparagraph:P104:end

% teaching: NEW-U06-GRAM-STRONG-IMPERATIVE-FORMS
% mdparagraph:P105:begin
The ordinary singular command uses the present stem without the infinitive ending. Do not import the changed vowel of the indicative second or third singular into it. The ordinary plural command matches the present indicative plural.
% mdparagraph:P105:end

% mdparagraph:P106:begin
Contracted verbs need a separate singular command in \textbf{-h}. Class I wrēon gives \textbf{wrīh}, with long ī; Class II tēon gives \textbf{tēoh}, with long ēo. Their plurals and infinitives stay contracted. In the table, drīfan and wrēon belong to Class I; bēodan and tēon to II; bindan, helpan, and drincan to III.
% mdparagraph:P106:end

% teaching: NEW-U06-IMP-drifan
% teaching: NEW-U06-NONFINITE-drifan
% teaching: NEW-U06-IMP-wreon
% teaching: NEW-U06-NONFINITE-wreon
% teaching: NEW-U06-IMP-beodan
% teaching: NEW-U06-NONFINITE-beodan
% teaching: NEW-U06-IMP-teon
% teaching: NEW-U06-NONFINITE-teon
% teaching: NEW-U06-IMP-bindan
% teaching: NEW-U06-NONFINITE-bindan
% teaching: NEW-U06-IMP-helpan
% teaching: NEW-U06-NONFINITE-helpan
% teaching: NEW-U06-IMP-drincan
% teaching: NEW-U06-NONFINITE-drincan
% table: NEW-U06-DISPLAY-STRONG-I-III
% mdtable:T10:begin
\par\begingroup
\fontsize{10.5}{13.2}\selectfont
\setlength{\LTleft}{0pt}\setlength{\LTright}{\fill}
\noindent
\begin{tabular}{@{}>{\raggedright\arraybackslash}p{17.31mm}>{\raggedright\arraybackslash}p{17.57mm}>{\raggedright\arraybackslash}p{23.8mm}>{\raggedright\arraybackslash}p{21.98mm}>{\raggedright\arraybackslash}p{32.1mm}@{}}
\toprule
\textbf{Verb} & \textbf{Command to one} & \textbf{Command to two or more} & \textbf{Plain infinitive} & \textbf{Inflected infinitive after tō} \\
\midrule
drīfan & drīf & drīfað & drīfan & drīfanne \\
wrēon & wrīh & wrēoð & wrēon & wrēonne \\
bēodan & bēod & bēodað & bēodan & bēodanne \\
tēon & tēoh & tēoð & tēon & tēonne \\
bindan & bind & bindað & bindan & bindanne \\
helpan & help & helpað & helpan & helpanne \\
drincan & drinc & drincað & drincan & drincanne \\
\bottomrule
\end{tabular}
\par\endgroup

% mdtable:T10:end
\reserve{7\baselineskip}
% mdparagraph:P107:begin
\subsection*{Strong IV–V}
% mdparagraph:P107:end

% mdparagraph:P108:begin
Beran, niman, and cuman are Class IV. The remaining verbs in this table are Class V. The ordinary singular \textbf{cum} retains u, unlike the changed indicative singular vowel.
% mdparagraph:P108:end

% mdparagraph:P109:begin
The strong j-present verbs have singular commands in \textbf{-e}, like the short weak-I pattern: \textbf{bide, lige, site}. Gemination is absent in these singulars. Licgan has g in lige, beside cg in licgað. Contracted sēon has \textbf{seoh}, with short eo, beside long ēo in sēoð.
% mdparagraph:P109:end

% teaching: NEW-U06-IMP-cuman
% teaching: NEW-U06-NONFINITE-cuman
% teaching: NEW-U06-IMP-beran
% teaching: NEW-U06-NONFINITE-beran
% teaching: NEW-U06-IMP-niman
% teaching: NEW-U06-NONFINITE-niman
% teaching: NEW-U06-IMP-giefan
% teaching: NEW-U06-NONFINITE-giefan
% teaching: NEW-U06-IMP-biddan
% teaching: NEW-U06-NONFINITE-biddan
% teaching: NEW-U06-IMP-licgan
% teaching: NEW-U06-NONFINITE-licgan
% teaching: NEW-U06-IMP-sittan
% teaching: NEW-U06-NONFINITE-sittan
% teaching: NEW-U06-IMP-etan
% teaching: NEW-U06-NONFINITE-etan
% teaching: NEW-U06-IMP-seon
% teaching: NEW-U06-NONFINITE-seon
% table: NEW-U06-DISPLAY-STRONG-IV-V
% mdtable:T11:begin
\par\begingroup
\fontsize{10.5}{13.2}\selectfont
\setlength{\LTleft}{0pt}\setlength{\LTright}{\fill}
\noindent
\begin{tabular}{@{}>{\raggedright\arraybackslash}p{17.31mm}>{\raggedright\arraybackslash}p{17.57mm}>{\raggedright\arraybackslash}p{23.8mm}>{\raggedright\arraybackslash}p{21.98mm}>{\raggedright\arraybackslash}p{32.1mm}@{}}
\toprule
\textbf{Verb} & \textbf{Command to one} & \textbf{Command to two or more} & \textbf{Plain infinitive} & \textbf{Inflected infinitive after tō} \\
\midrule
cuman & cum & cumað & cuman & cumanne \\
beran & ber & berað & beran & beranne \\
niman & nim & nimað & niman & nimanne \\
giefan & gief & giefað & giefan & giefanne \\
biddan & bide & biddað & biddan & biddanne \\
licgan & lige & licgað & licgan & licganne \\
sittan & site & sittað & sittan & sittanne \\
etan & et & etað & etan & etanne \\
sēon & seoh & sēoð & sēon & sēonne \\
\bottomrule
\end{tabular}
\par\endgroup

% mdtable:T11:end
\reserve{7\baselineskip}
% mdparagraph:P110:begin
\subsection*{Strong VI–VII}
% mdparagraph:P110:end

% mdparagraph:P111:begin
Faran, standan, slēan, hębban, scieppan, and swęrian are Class VI. Hātan, healdan, weaxan, lǣtan, and fōn are Class VII.
% mdparagraph:P111:end

% mdparagraph:P112:begin
The j-present singulars are \textbf{hęfe, sciepe, swęre}. Hębban changes bb to f in the singular command. Swęrian has swęre beside swęriað. Contracted slēan gives \textbf{sleah}, with short ea; fōn gives \textbf{fōh}, with long ō. Keep these quantities distinct from those in the corresponding infinitives.
% mdparagraph:P112:end

% teaching: NEW-U06-IMP-hatan
% teaching: NEW-U06-NONFINITE-hatan
% teaching: NEW-U06-IMP-standan
% teaching: NEW-U06-NONFINITE-standan
% teaching: NEW-U06-IMP-faran
% teaching: NEW-U06-NONFINITE-faran
% teaching: NEW-U06-IMP-slean
% teaching: NEW-U06-NONFINITE-slean
% teaching: NEW-U06-IMP-hebban
% teaching: NEW-U06-NONFINITE-hebban
% teaching: NEW-U06-IMP-scieppan
% teaching: NEW-U06-NONFINITE-scieppan
% teaching: NEW-U06-IMP-swerian
% teaching: NEW-U06-NONFINITE-swerian
% teaching: NEW-U06-IMP-healdan
% teaching: NEW-U06-NONFINITE-healdan
% teaching: NEW-U06-IMP-weaxan
% teaching: NEW-U06-NONFINITE-weaxan
% teaching: NEW-U06-IMP-laetan
% teaching: NEW-U06-NONFINITE-laetan
% teaching: NEW-U06-IMP-fon
% teaching: NEW-U06-NONFINITE-fon
% table: NEW-U06-DISPLAY-STRONG-VI-VII
% mdtable:T12:begin
\par\begingroup
\fontsize{10.5}{13.2}\selectfont
\setlength{\LTleft}{0pt}\setlength{\LTright}{\fill}
\noindent
\begin{tabular}{@{}>{\raggedright\arraybackslash}p{17.16mm}>{\raggedright\arraybackslash}p{17.42mm}>{\raggedright\arraybackslash}p{24.54mm}>{\raggedright\arraybackslash}p{21.8mm}>{\raggedright\arraybackslash}p{31.83mm}@{}}
\toprule
\textbf{Verb} & \textbf{Command to one} & \textbf{Command to two or more} & \textbf{Plain infinitive} & \textbf{Inflected infinitive after tō} \\
\midrule
hātan & hāt & hātað & hātan & hātanne \\
standan & stand & standað & standan & standanne \\
faran & far & farað & faran & faranne \\
slēan & sleah & slēað & slēan & slēanne \\
hębban & hęfe & hębbað & hębban & hębbanne \\
scieppan & sciepe & scieppað & scieppan & scieppanne \\
swęrian & swęre & swęriað & swęrian & swęrianne \\
healdan & heald & healdað & healdan & healdanne \\
weaxan & weax & weaxað & weaxan & weaxanne \\
lǣtan & lǣt & lǣtað & lǣtan & lǣtanne \\
fōn & fōh & fōð & fōn & fōnne \\
\bottomrule
\end{tabular}
\par\endgroup

% mdtable:T12:end
% mdparagraph:P113:begin
\textbf{Sources for the strong forms.} Cook §§93–95, pp.54–55 /\allowbreak{} PDF pp.70–71; §§100–101, pp.57–59 /\allowbreak{} PDF pp.73–75; §§106–107, pp.64–65 /\allowbreak{} PDF pp.80–81. Wright §§478 and 480, pp.250–251 /\allowbreak{} PDF pp.270–271, corroborates imperative endings and inflected infinitives. Fulk §127, pp.93–94 /\allowbreak{} PDF pp.105–106, corroborates the j-present singulars and their consonant alternations. Wrīh applies Cook’s stated Class I contracted pattern to wrēon; the model-derived form is not presented as a separate prose quotation.
% mdparagraph:P113:end

\reserve{7\baselineskip}
% mdparagraph:P114:begin
\subsection*{Dōn, gān, wesan, and bēon}
% mdparagraph:P114:end

% teaching: NEW-U06-GRAM-ANOMALOUS-IMPERATIVE-FORMS
% mdparagraph:P115:begin
Learn \textbf{dō/\allowbreak{}dōð}, \textbf{gā/\allowbreak{}gāð}, \textbf{wes/\allowbreak{}wesað}, and \textbf{bēo/\allowbreak{}bēoð} as complete pairs. Their contracted infinitives keep the contracted stem before -nne. Wesan has the ordinary -anne form.
% mdparagraph:P115:end

% mdparagraph:P116:begin
These second-person commands are distinct from an indicative or a supplied third-person wish that may share a spelling. Neither willan nor the preterite-present verbs receives a regular command pair by analogy with them.
% mdparagraph:P116:end

% teaching: NEW-U06-IMP-wesan
% teaching: NEW-U06-NONFINITE-wesan
% teaching: NEW-U06-IMP-don
% teaching: NEW-U06-NONFINITE-don
% teaching: NEW-U06-IMP-gan
% teaching: NEW-U06-NONFINITE-gan
% teaching: NEW-U06-IMP-wesan-beon
% teaching: NEW-U06-NONFINITE-wesan-beon
% table: NEW-U06-DISPLAY-ANOMALOUS
% mdtable:T13:begin
\par\begingroup
\fontsize{10.5}{13.2}\selectfont
\setlength{\LTleft}{0pt}\setlength{\LTright}{\fill}
\noindent
\begin{tabular}{@{}>{\raggedright\arraybackslash}p{17.31mm}>{\raggedright\arraybackslash}p{17.57mm}>{\raggedright\arraybackslash}p{23.8mm}>{\raggedright\arraybackslash}p{21.98mm}>{\raggedright\arraybackslash}p{32.1mm}@{}}
\toprule
\textbf{Verb} & \textbf{Command to one} & \textbf{Command to two or more} & \textbf{Plain infinitive} & \textbf{Inflected infinitive after tō} \\
\midrule
wesan & wes & wesað & wesan & wesanne \\
dōn & dō & dōð & dōn & dōnne \\
gān & gā & gāð & gān & gānne \\
bēon & bēo & bēoð & bēon & bēonne \\
\bottomrule
\end{tabular}
\par\endgroup

% mdtable:T13:end
% example: NEW-U06-EXAMPLE-08
% mdparagraph:P117:begin
Gā. Dō ðæt weorc wel.
% mdparagraph:P117:end

% mdparagraph:P118:begin
Go. Do the work well, one addressee.
% mdparagraph:P118:end

% mdparagraph:P119:begin
Some verbs remain form models for uses not yet learned. A full command or infinitive row does not by itself teach a new object, a causative complement, or a naming construction. Keep the earlier limits on such models as gierwan and hycgan, and on the complement uses of hātan, lǣtan, bēodan, and biddan.
% mdparagraph:P119:end

% mdparagraph:P120:begin
\textbf{Reference variant.} Cook records a shortened plural \textbf{-e} immediately before a following subject pronoun. Use the full plural forms above when the subject is omitted or precedes the verb. The positional variant is supplied for reference here, not an additional free choice of ending or a first-/\allowbreak{}third-person command system.
% mdparagraph:P120:end

% mdparagraph:P121:begin
\textbf{Sources.} Cook §§138–141, pp.78–80 /\allowbreak{} PDF pp.94–96; Fulk §134, p.102 /\allowbreak{} PDF p.114. The postposed-pronoun variant is Cook §95, p.55 /\allowbreak{} PDF p.71. Inflected -anne/\allowbreak{}-enne is also described by Wright §480, p.251 /\allowbreak{} PDF p.271, and Fulk §58, p.38 /\allowbreak{} PDF p.50.
% mdparagraph:P121:end

\reserve{7\baselineskip}
% mdparagraph:P122:begin
\subsection*{Objects, negative commands, and quoted speech}
% mdparagraph:P122:end

% mdparagraph:P123:begin
Commands retain the verb’s familiar government. Giving still has a dative recipient and an accusative thing; helpan still takes the dative person helped.
% mdparagraph:P123:end

% example: NEW-U06-EXAMPLE-05
% mdparagraph:P124:begin
Gief mē ðæt wæter. Helpað mē.
% mdparagraph:P124:end

% mdparagraph:P125:begin
Give me the water, one addressee. Help me, several addressees.
% mdparagraph:P125:end

% teaching: NEW-U06-GRAM-IMPERATIVE-NEGATIVE
% mdparagraph:P126:begin
Put \textbf{ne} directly before the imperative for the simple negative command practised here.
% mdparagraph:P126:end

% example: NEW-U06-EXAMPLE-06
% mdparagraph:P127:begin
Ne drinc ðæt wæter.
% mdparagraph:P127:end

% mdparagraph:P128:begin
Do not drink the water, one addressee.
% mdparagraph:P128:end

% teaching: NEW-U06-GRAM-DIRECT-QUOTATION
% mdparagraph:P129:begin
A quotation introduced by \textbf{cweðan} may now contain these direct commands. Within the quotation, “I,” “me,” and “you” belong to the speaker’s and addressee’s situation. The quotation is not an indirect command and does not require a new dependent mood.
% mdparagraph:P129:end

% mdparagraph:P130:begin
\textbf{Sources.} Cook §198, p.105 /\allowbreak{} PDF p.121, gives commands with or without an expressed second-person subject. The negative imperative is directly evidenced in Reader XII, p.183 /\allowbreak{} PDF p.199, editorial line 7, and quoted in §§156h and 159a, pp.92–93 /\allowbreak{} PDF pp.108–109. These citations support the construction; their additional vocabulary is not needed for the exercises.
% mdparagraph:P130:end

\reserve{7\baselineskip}
% mdparagraph:P131:begin
\section*{5. Practice in clauses and forms}
% mdparagraph:P131:end

% mdparagraph:P132:begin
All exercises are new teaching material. Unless an exercise asks for a supplied reference, use the forms and constructions learned in this and earlier units. Preserve the stated modal sense. When a particular course form or construction is requested, distinguish that target from another valid expression.
% mdparagraph:P132:end

\reserve{7\baselineskip}
% mdparagraph:P133:begin
\subsection*{Practice M. Meanings and infinitive complements}
% mdparagraph:P133:end

\reserve{5\baselineskip}
\par\begin{minipage}{\linewidth}\raggedright
% mdparagraph:P134:begin
\textbf{M1.}
% mdparagraph:P134:end

% item: NEW-U06-M01
% mdparagraph:P135:begin
Translate “I know how to swim. We know how to sing.” Use cunnan for learned ability.
% mdparagraph:P135:end

\end{minipage}\par\addvspace{4pt}

\par\begin{minipage}{\linewidth}\raggedright
% mdparagraph:P136:begin
\textbf{M2.}
% mdparagraph:P136:end

% item: NEW-U06-M02
% mdparagraph:P137:begin
Translate “He can carry the water. They can help the man.” Use mugan for ability.
% mdparagraph:P137:end

\end{minipage}\par\addvspace{4pt}

\par\begin{minipage}{\linewidth}\raggedright
% mdparagraph:P138:begin
\textbf{M3.}
% mdparagraph:P138:end

% item: NEW-U06-M03
% mdparagraph:P139:begin
Translate “I need to go. We needed to go.” Use ðurfan.
% mdparagraph:P139:end

\end{minipage}\par\addvspace{4pt}

\par\begin{minipage}{\linewidth}\raggedright
% mdparagraph:P140:begin
\textbf{M4.}
% mdparagraph:P140:end

% item: NEW-U06-M04
% mdparagraph:P141:begin
Translate “You dare to go. You dared to go.” Address one person and use durran.
% mdparagraph:P141:end

\end{minipage}\par\addvspace{4pt}

\par\begin{minipage}{\linewidth}\raggedright
% mdparagraph:P142:begin
\textbf{M5.}
% mdparagraph:P142:end

% item: NEW-U06-M05
% mdparagraph:P143:begin
Translate “We must do the work. They had to do the work.” Use sculan for obligation.
% mdparagraph:P143:end

\end{minipage}\par\addvspace{4pt}

\par\begin{minipage}{\linewidth}\raggedright
% mdparagraph:P144:begin
\textbf{M6.}
% mdparagraph:P144:end

% item: NEW-U06-M06
% mdparagraph:P145:begin
Translate “You are allowed to eat the food. You were allowed to eat the food.” Address one person and use mōtan.
% mdparagraph:P145:end

\end{minipage}\par\addvspace{4pt}

\par\begin{minipage}{\linewidth}\raggedright
% mdparagraph:P146:begin
\textbf{M7.}
% mdparagraph:P146:end

% item: NEW-U06-M07
% mdparagraph:P147:begin
Translate “He wants to give the birds water. They wanted to give the birds water.” Use willan for desire.
% mdparagraph:P147:end

\end{minipage}\par\addvspace{4pt}

\par\begin{minipage}{\linewidth}\raggedright
% mdparagraph:P148:begin
\textbf{M8.}
% mdparagraph:P148:end

% item: NEW-U06-M08
% mdparagraph:P149:begin
Translate “I did not dare to go. We cannot carry the ship.” Use the taught simple ne construction.
% mdparagraph:P149:end

\end{minipage}\par\addvspace{4pt}

\par\begin{minipage}{\linewidth}\raggedright
% mdparagraph:P150:begin
\textbf{M9.}
% mdparagraph:P150:end

% item: NEW-U06-M09
% mdparagraph:P151:begin
Translate “We two know how to swim. You two must help us two.” Keep the two-person references.
% mdparagraph:P151:end

\end{minipage}\par\addvspace{4pt}

\par\begin{minipage}{\linewidth}\raggedright
% mdparagraph:P152:begin
\textbf{M10.}
% mdparagraph:P152:end

% item: NEW-U06-M10
% mdparagraph:P153:begin
Translate “He wanted to go. He went.” Explain why the first sentence alone does not assert that he went.
% mdparagraph:P153:end

\end{minipage}\par\addvspace{4pt}

\par\begin{minipage}{\linewidth}\raggedright
% mdparagraph:P154:begin
\textbf{M11.}
% mdparagraph:P154:end

% item: NEW-U06-M11
% mdparagraph:P155:begin
Translate the printed pairs Ic cann swimman and Ic mæg swimman using the unit’s stated difference between learned ability and present capability. Identify the finite verb and the infinitive in each. Do not claim that one English word always represents only one Old English modal.
% mdparagraph:P155:end

\end{minipage}\par\addvspace{4pt}

\par\begin{minipage}{\linewidth}\raggedright
% mdparagraph:P156:begin
\textbf{M12.}
% mdparagraph:P156:end

% item: NEW-U06-M12
% mdparagraph:P157:begin
Translate “The ship will be ready.” Use the bēo-series indicative to describe the expected later state. Then contrast the earlier eom-series sentence “The ship is ready.” Do not infer that every bēo-series occurrence must be future.
% mdparagraph:P157:end

\end{minipage}\par\addvspace{4pt}

\par\begin{minipage}{\linewidth}\raggedright
% mdparagraph:P158:begin
\textbf{M13.}
% mdparagraph:P158:end

% item: NEW-U06-M13
% mdparagraph:P159:begin
The “things” here are matters already discussed. Translate “I know those things. He knew those things.” Use witan. Explain how this knowledge construction differs from cunnan followed by an infinitive.
% mdparagraph:P159:end

\end{minipage}\par\addvspace{4pt}

\reserve{7\baselineskip}
% mdparagraph:P160:begin
\subsection*{Practice I. Complete command pairs and instructions}
% mdparagraph:P160:end

\reserve{5\baselineskip}
\par\begin{minipage}{\linewidth}\raggedright
% mdparagraph:P161:begin
\textbf{I1.}
% mdparagraph:P161:end

% item: NEW-U06-I01
% mdparagraph:P162:begin
Give the second-person singular and plural imperatives of hīeran, nęrian and gefręmman. Label both numbers.
% mdparagraph:P162:end

\end{minipage}\par\addvspace{4pt}

\par\begin{minipage}{\linewidth}\raggedright
% mdparagraph:P163:begin
\textbf{I2.}
% mdparagraph:P163:end

% item: NEW-U06-I02
% mdparagraph:P164:begin
Give the second-person singular and plural imperatives of cyssan, gefyllan, sęttan, gierwan and sęllan. Label both numbers.
% mdparagraph:P164:end

\end{minipage}\par\addvspace{4pt}

\par\begin{minipage}{\linewidth}\raggedright
% mdparagraph:P165:begin
\textbf{I3.}
% mdparagraph:P165:end

% item: NEW-U06-I03
% mdparagraph:P166:begin
Give the second-person singular and plural imperatives of lufian, habban and libban. Label both numbers.
% mdparagraph:P166:end

\end{minipage}\par\addvspace{4pt}

\par\begin{minipage}{\linewidth}\raggedright
% mdparagraph:P167:begin
\textbf{I4.}
% mdparagraph:P167:end

% item: NEW-U06-I04
% mdparagraph:P168:begin
Give the second-person singular and plural imperatives of sęcgan and hycgan. Label both numbers.
% mdparagraph:P168:end

\end{minipage}\par\addvspace{4pt}

\par\begin{minipage}{\linewidth}\raggedright
% mdparagraph:P169:begin
\textbf{I5.}
% mdparagraph:P169:end

% item: NEW-U06-I05
% mdparagraph:P170:begin
Give the second-person singular and plural imperatives of drīfan, bēodan, bindan, helpan and drincan. Label both numbers.
% mdparagraph:P170:end

\end{minipage}\par\addvspace{4pt}

\par\begin{minipage}{\linewidth}\raggedright
% mdparagraph:P171:begin
\textbf{I6.}
% mdparagraph:P171:end

% item: NEW-U06-I06
% mdparagraph:P172:begin
Give the second-person singular and plural imperatives of wrēon and tēon. Label both numbers.
% mdparagraph:P172:end

\end{minipage}\par\addvspace{4pt}

\par\begin{minipage}{\linewidth}\raggedright
% mdparagraph:P173:begin
\textbf{I7.}
% mdparagraph:P173:end

% item: NEW-U06-I07
% mdparagraph:P174:begin
Give the second-person singular and plural imperatives of beran, niman and cuman. Label both numbers.
% mdparagraph:P174:end

\end{minipage}\par\addvspace{4pt}

\par\begin{minipage}{\linewidth}\raggedright
% mdparagraph:P175:begin
\textbf{I8.}
% mdparagraph:P175:end

% item: NEW-U06-I08
% mdparagraph:P176:begin
Give the second-person singular and plural imperatives of giefan, biddan, sittan, licgan and etan. Label both numbers.
% mdparagraph:P176:end

\end{minipage}\par\addvspace{4pt}

\par\begin{minipage}{\linewidth}\raggedright
% mdparagraph:P177:begin
\textbf{I9.}
% mdparagraph:P177:end

% item: NEW-U06-I09
% mdparagraph:P178:begin
Give the second-person singular and plural imperatives of sēon, slēan and fōn. Label both numbers.
% mdparagraph:P178:end

\end{minipage}\par\addvspace{4pt}

\par\begin{minipage}{\linewidth}\raggedright
% mdparagraph:P179:begin
\textbf{I10.}
% mdparagraph:P179:end

% item: NEW-U06-I10
% mdparagraph:P180:begin
Give the second-person singular and plural imperatives of hębban, scieppan and swęrian. Label both numbers.
% mdparagraph:P180:end

\end{minipage}\par\addvspace{4pt}

\par\begin{minipage}{\linewidth}\raggedright
% mdparagraph:P181:begin
\textbf{I11.}
% mdparagraph:P181:end

% item: NEW-U06-I11
% mdparagraph:P182:begin
Give the second-person singular and plural imperatives of faran, standan and healdan. Label both numbers.
% mdparagraph:P182:end

\end{minipage}\par\addvspace{4pt}

\par\begin{minipage}{\linewidth}\raggedright
% mdparagraph:P183:begin
\textbf{I12.}
% mdparagraph:P183:end

% item: NEW-U06-I12
% mdparagraph:P184:begin
Give the second-person singular and plural imperatives of weaxan, hātan and lǣtan. Label both numbers.
% mdparagraph:P184:end

\end{minipage}\par\addvspace{4pt}

\par\begin{minipage}{\linewidth}\raggedright
% mdparagraph:P185:begin
\textbf{I13.}
% mdparagraph:P185:end

% item: NEW-U06-I13
% mdparagraph:P186:begin
Give the second-person singular and plural imperatives of dōn, gān and both command series of wesan/\allowbreak{}bēon. Label both numbers.
% mdparagraph:P186:end

\end{minipage}\par\addvspace{4pt}

\par\begin{minipage}{\linewidth}\raggedright
% mdparagraph:P187:begin
\textbf{I14.}
% mdparagraph:P187:end

% item: NEW-U06-I14
% mdparagraph:P188:begin
Translate “Take the food and carry it.” First address one person, then several. Identify what “it” refers to and preserve its grammatical gender.
% mdparagraph:P188:end

\end{minipage}\par\addvspace{4pt}

\par\begin{minipage}{\linewidth}\raggedright
% mdparagraph:P189:begin
\textbf{I15.}
% mdparagraph:P189:end

% item: NEW-U06-I15
% mdparagraph:P190:begin
Translate “Help me. Give me the water.” First address one person, then exactly two people. Explain the verb number used for the two addressees.
% mdparagraph:P190:end

\end{minipage}\par\addvspace{4pt}

\par\begin{minipage}{\linewidth}\raggedright
% mdparagraph:P191:begin
\textbf{I16.}
% mdparagraph:P191:end

% item: NEW-U06-I16
% mdparagraph:P192:begin
Translate “Do not drink the water. Sit and eat.” Address several people.
% mdparagraph:P192:end

\end{minipage}\par\addvspace{4pt}

\par\begin{minipage}{\linewidth}\raggedright
% mdparagraph:P193:begin
\textbf{I17.}
% mdparagraph:P193:end

% item: NEW-U06-I17
% mdparagraph:P194:begin
Translate “The man said: ‘Go and do the work well.’” The man addresses one companion.
% mdparagraph:P194:end

\end{minipage}\par\addvspace{4pt}

\reserve{7\baselineskip}
% mdparagraph:P195:begin
\subsection*{Practice N. Infinitives and purpose}
% mdparagraph:P195:end

\reserve{5\baselineskip}
\par\begin{minipage}{\linewidth}\raggedright
% mdparagraph:P196:begin
\textbf{N1.}
% mdparagraph:P196:end

% item: NEW-U06-N01
% mdparagraph:P197:begin
Give the plain infinitive and the selected inflected infinitive after tō for bindan, beran and sęllan. Use the course’s -anne forms where the table selects them.
% mdparagraph:P197:end

\end{minipage}\par\addvspace{4pt}

\par\begin{minipage}{\linewidth}\raggedright
% mdparagraph:P198:begin
\textbf{N2.}
% mdparagraph:P198:end

% item: NEW-U06-N02
% mdparagraph:P199:begin
Give the plain infinitive and the selected inflected infinitive after tō for hīeran, cyssan, nęrian, gefręmman and gierwan. Use the course’s -anne forms where the table selects them.
% mdparagraph:P199:end

\end{minipage}\par\addvspace{4pt}

\par\begin{minipage}{\linewidth}\raggedright
% mdparagraph:P200:begin
\textbf{N3.}
% mdparagraph:P200:end

% item: NEW-U06-N03
% mdparagraph:P201:begin
Give the plain infinitive and the selected inflected infinitive after tō for lufian, habban, libban, sęcgan and hycgan. Use the course’s -anne forms where the table selects them.
% mdparagraph:P201:end

\end{minipage}\par\addvspace{4pt}

\par\begin{minipage}{\linewidth}\raggedright
% mdparagraph:P202:begin
\textbf{N4.}
% mdparagraph:P202:end

% item: NEW-U06-N04
% mdparagraph:P203:begin
Give the plain infinitive and the selected inflected infinitive after tō for wrēon, tēon, sēon, slēan and fōn. Use the course’s -anne forms where the table selects them.
% mdparagraph:P203:end

\end{minipage}\par\addvspace{4pt}

\par\begin{minipage}{\linewidth}\raggedright
% mdparagraph:P204:begin
\textbf{N5.}
% mdparagraph:P204:end

% item: NEW-U06-N05
% mdparagraph:P205:begin
Give the plain infinitive and the selected inflected infinitive after tō for gān, dōn and both wesan/\allowbreak{}bēon series. Use the course’s -anne forms where the table selects them.
% mdparagraph:P205:end

\end{minipage}\par\addvspace{4pt}

\par\begin{minipage}{\linewidth}\raggedright
% mdparagraph:P206:begin
\textbf{N6.}
% mdparagraph:P206:end

% item: NEW-U06-N06
% mdparagraph:P207:begin
Translate “He wants to eat. He came to eat.” Use the taught plain-infinitive complement in the first sentence and the taught purpose construction in the second.
% mdparagraph:P207:end

\end{minipage}\par\addvspace{4pt}

\par\begin{minipage}{\linewidth}\raggedright
% mdparagraph:P208:begin
\textbf{N7.}
% mdparagraph:P208:end

% item: NEW-U06-N07
% mdparagraph:P209:begin
Translate “The companions came to help the man.” Use the inflected infinitive for their purpose. Identify the person doing the helping and the case of the person helped.
% mdparagraph:P209:end

\end{minipage}\par\addvspace{4pt}

\par\begin{minipage}{\linewidth}\raggedright
% mdparagraph:P210:begin
\textbf{N8.}
% mdparagraph:P210:end

% item: NEW-U06-N08
% mdparagraph:P211:begin
Translate “The man came to give food to the birds.” Use the inflected infinitive for the purpose.
% mdparagraph:P211:end

\end{minipage}\par\addvspace{4pt}

\par\begin{minipage}{\linewidth}\raggedright
% mdparagraph:P212:begin
\textbf{N9.}
% mdparagraph:P212:end

% item: NEW-U06-N09
% mdparagraph:P213:begin
In Hīe sculon ðæt weorc dōn and Hīe cōmon tō dōnne ðæt weorc, identify each finite verb, each infinitive, and the construction that selects the infinitive form. Translate both sentences.
% mdparagraph:P213:end

\end{minipage}\par\addvspace{4pt}

\par\begin{minipage}{\linewidth}\raggedright
% mdparagraph:P214:begin
\textbf{N10.}
% mdparagraph:P214:end

% item: NEW-U06-N10
% mdparagraph:P215:begin
Complete the taught modal construction “We can carry the food” using mugan and beran. Then make the subject “I”. Explain which verb changes for person and number.
% mdparagraph:P215:end

\end{minipage}\par\addvspace{4pt}

\reserve{7\baselineskip}
% mdparagraph:P216:begin
\section*{6. Food and water for the journey}
% mdparagraph:P216:end

% mdparagraph:P217:begin
The following connected account is new teaching material. Read its statements and instructions together. A command requests an action; the later narrative reports what occurred.
% mdparagraph:P217:end

% composition: NEW-U06-TEXT-A
% mdparagraph:P218:begin
\par\begin{minipage}{\linewidth}\raggedright
\readingtext{Se mann cwæð: ‘Ic wille faran. Mīn scip is gearu. Cumað and helpað mē.’ Twēgen gefēran cōmon tō helpanne. Se mann cwæð: ‘Nimað ðone męte and berað hine. Giefað ðǣm fuglum ðæt wæter.’ Se gefēra cwæð: ‘Ic mæg ðone męte beran. Wē sculon ðæt weorc gefręmman.’ Ðā gefēran bǣron ðone męte and gēafon ðǣm fuglum ðæt wæter. Se mann cwæð: ‘Sittað and etað.’ Ðā gefēran ǣton. Hīe fōron.}
\end{minipage}\par
% mdparagraph:P218:end

\reserve{7\baselineskip}
% mdparagraph:P219:begin
\subsection*{Practice C. The connected account}
% mdparagraph:P219:end

\reserve{5\baselineskip}
\par\begin{minipage}{\linewidth}\raggedright
% mdparagraph:P220:begin
\textbf{C1.}
% mdparagraph:P220:end

% item: NEW-U06-C01
% mdparagraph:P221:begin
Read “Food and water for the journey”. Who asks for help, how many companions come, and what is their stated purpose? Answer in English.
% mdparagraph:P221:end

\end{minipage}\par\addvspace{4pt}

\par\begin{minipage}{\linewidth}\raggedright
% mdparagraph:P222:begin
\textbf{C2.}
% mdparagraph:P222:end

% item: NEW-U06-C02
% mdparagraph:P223:begin
List every imperative in the controlled text and give its addressee number. For each, give the singular command form taught in the tables.
% mdparagraph:P223:end

\end{minipage}\par\addvspace{4pt}

\par\begin{minipage}{\linewidth}\raggedright
% mdparagraph:P224:begin
\textbf{C3.}
% mdparagraph:P224:end

% item: NEW-U06-C03
% mdparagraph:P225:begin
In the controlled text, identify the antecedent and case of hine, the recipient of the water, and what the companions subsequently eat. Distinguish what is explicit from what the context suggests.
% mdparagraph:P225:end

\end{minipage}\par\addvspace{4pt}

\par\begin{minipage}{\linewidth}\raggedright
% mdparagraph:P226:begin
\textbf{C4.}
% mdparagraph:P226:end

% item: NEW-U06-C04
% mdparagraph:P227:begin
Identify wille, mæg and sculon in the controlled text by person, number, tense and stated sense. Give the infinitive governed by each.
% mdparagraph:P227:end

\end{minipage}\par\addvspace{4pt}

\par\begin{minipage}{\linewidth}\raggedright
% mdparagraph:P228:begin
\textbf{C5.}
% mdparagraph:P228:end

% item: NEW-U06-C05
% mdparagraph:P229:begin
Translate the whole controlled text. Preserve the distinctions among a wish to travel, ability to carry, obligation, commands and completed past actions.
% mdparagraph:P229:end

\end{minipage}\par\addvspace{4pt}

\par\begin{minipage}{\linewidth}\raggedright
% mdparagraph:P230:begin
\textbf{C6.}
% mdparagraph:P230:end

% item: NEW-U06-C06
% mdparagraph:P231:begin
Rewrite the man’s first quoted commands, Cumað and helpað mē, as commands to one companion. Separately, put the first sentence of the original narrative into the present, keeping its original quotation unchanged.
% mdparagraph:P231:end

\end{minipage}\par\addvspace{4pt}

\reserve{7\baselineskip}
% mdparagraph:P232:begin
\section*{7. The fisherman’s gift and first direction}
% mdparagraph:P232:end

% mdparagraph:P233:begin
\textbf{Editorial context.} Apollonius survived the earlier storm, encountered an old fisherman, and asked for help.
% mdparagraph:P233:end

% mdparagraph:P234:begin
This excerpt continues with the fisherman’s care and gift, ending after his first complete spoken direction. It is from Cook’s Reader XII, \textbf{Apollonius}, printed p.167 /\allowbreak{} one-based PDF p.183, editorial prose lines 3–10. The edition has printed marginal anchors at 5 and 10; the intervening line numbers are editorial. The speech continues in the source. A closing quotation mark is added at this excerpt’s boundary; the Old English wording and spellings are retained.
% mdparagraph:P234:end

% source: FB-XII-APOLLONIUS-FISHERMAN-GIFT-01
% mdparagraph:P235:begin
\par\begin{minipage}{\linewidth}\raggedright
\readingtext{Ðā, sōna swā se fiscere geseah ðæt se geonga mann æt his fōtum læg, hē mid mildheortnesse hine ūpāhōf, and lǣdde hine mid him tō his hūse, and ðā ēstas him beforan lęgde ðe hē him tō bēodanne hæfde. Ðā gīet hē wolde, be his mihte, māran ārfæstnesse him gecȳðan; tōslāt ðā his wǣfels on twā, and sealde Apollonie ðone healfan dǣl, ðus cweðende: “Nim ðæt ic ðē tō sęllanne hæbbe, and gā intō ðǣre ceastre.”}
\end{minipage}\par
% mdparagraph:P235:end

\reserve{7\baselineskip}
% mdparagraph:P236:begin
\subsection*{Local help}
% mdparagraph:P236:end

% mdparagraph:P237:begin
These ten groups of help remain available during the reading exercises. Their unfamiliar forms and constructions are supplied for this passage; they are not independent-recall vocabulary or general construction patterns. The learned finite verbs, commands, and infinitive morphology can still be identified within the supported clauses.
% mdparagraph:P237:end

\reserve{5\baselineskip}
% mdparagraph:P238:begin
\textbf{1. Ðā, sōna swā se fiscere geseah ðæt se geonga mann æt his fōtum læg,}
% mdparagraph:P238:end

% support: U06-SRC-SUP-OPENING
% teaching: NEW-U06-SUPFORM-tha-adverb-invariant
% teaching: NEW-U06-SUPSENSE-tha-adverb-1
% teaching: NEW-U06-SUPFORM-sona-invariant-source
% teaching: NEW-U06-SUPSENSE-sona-as-soon-source
% teaching: NEW-U06-SUPFORM-swa-invariant
% teaching: NEW-U06-SUPSENSE-swa-as-correlative-source
% teaching: NEW-U06-SUPFORM-fiscere-nom-sg-source
% teaching: NEW-U06-SUPSENSE-fiscere-fisherman-source
% teaching: NEW-U06-SUPFORM-thaet-conjunction-invariant
% teaching: NEW-U06-SUPSENSE-thaet-conjunction-1
% teaching: NEW-U06-SUPFORM-geong-nom-m-sg-weak-source
% teaching: NEW-U06-SUPSENSE-geong-young-source
% teaching: NEW-U06-SUPFORM-aet-invariant-source
% teaching: NEW-U06-SUPSENSE-aet-at-source
% teaching: NEW-U06-SUPFORM-fot-dat-pl-source
% teaching: NEW-U06-SUPSENSE-fot-foot-source
% teaching: NEW-U06-GRAM-U06-SOURCE-TEMPORAL-CONTENT
% teaching: NEW-U06-GRAM-U06-SOURCE-PREPOSITIONS
% mdparagraph:P239:begin
Then, as soon as the fisherman saw that the young man was lying at his feet.
% mdparagraph:P239:end

% mdparagraph:P240:begin
\textbf{Ðā} means “then”; \textbf{sōna swā}, “as soon as.” Fiscere is nominative singular, “fisherman”; geonga is the weak masculine nominative singular adjective, “young.” The temporal clause contains the whole content clause introduced by ðæt, “that.” Its finite verb læg means “lay; was lying.” In æt his fōtum, “at his feet,” fōtum is dative plural. The feet are the fisherman’s. Both nested clause relations and the prepositional phrase are supplied as a whole.
% mdparagraph:P240:end

\reserve{5\baselineskip}
% mdparagraph:P241:begin
\textbf{2. hē mid mildheortnesse hine ūpāhōf,}
% mdparagraph:P241:end

% support: U06-SRC-SUP-COMPASSION
% teaching: NEW-U06-SUPFORM-mid-invariant-source
% teaching: NEW-U06-SUPSENSE-mid-manner-source
% teaching: NEW-U06-SUPFORM-mildheortnes-dat-sg-source
% teaching: NEW-U06-SUPSENSE-mildheortnes-compassion-source
% teaching: NEW-U06-SUPFORM-up-ahebban-pret-3-sg-source
% teaching: NEW-U06-SUPSENSE-up-ahebban-lift-source
% mdparagraph:P242:begin
He compassionately lifted him up.
% mdparagraph:P242:end

% mdparagraph:P243:begin
\textbf{Mid mildheortnesse}, “with compassion,” is a supplied manner phrase with dative mildheortnesse. \textbf{Ūpāhōf} means “lifted up”: it is the supplied third-singular preterite of a prefixed verb, not the unprefixed hębban form. Hē is the fisherman; hine is Apollonius.
% mdparagraph:P243:end

\reserve{5\baselineskip}
% mdparagraph:P244:begin
\textbf{3. and lǣdde hine mid him tō his hūse,}
% mdparagraph:P244:end

% support: U06-SRC-SUP-HOUSE
% teaching: NEW-U06-SUPSENSE-mid-accompaniment-source
% teaching: NEW-U06-SUPFORM-to-preposition-dative-invariant
% teaching: NEW-U06-SUPSENSE-to-destination
% teaching: NEW-U06-SUPFORM-hus-dat-sg-source
% teaching: NEW-U06-SUPSENSE-hus-house-source
% mdparagraph:P245:begin
And led him with him to his house.
% mdparagraph:P245:end

% mdparagraph:P246:begin
\textbf{Mid him} means “with him,” here accompanying the fisherman. \textbf{Tō his hūse} means “to his house”; hūse is the supplied dative singular of hūs, “house.” In this clause hine is Apollonius, while him and his refer contextually to the fisherman. The accompaniment and destination phrases are supplied locally. Tō before the noun phrase is distinct from tō before an inflected infinitive.
% mdparagraph:P246:end

\reserve{5\baselineskip}
% mdparagraph:P247:begin
\textbf{4. and ðā ēstas him beforan lęgde ðe hē him tō bēodanne hæfde.}
% mdparagraph:P247:end

% support: U06-SRC-SUP-PROVISIONS
% teaching: NEW-U06-SUPFORM-est-acc-m-pl-source
% teaching: NEW-U06-SUPSENSE-est-food-source
% teaching: NEW-U06-SUPFORM-beforan-invariant-source
% teaching: NEW-U06-SUPSENSE-beforan-before-source
% teaching: NEW-U06-SUPFORM-lecgean-pret-3-sg-source
% teaching: NEW-U06-SUPSENSE-lecgean-place-source
% teaching: NEW-U06-SUPFORM-the-relative-invariant
% teaching: NEW-U06-SUPSENSE-the-relative-which
% teaching: NEW-U06-GRAM-U06-SOURCE-AVAILABILITY
% teaching: NEW-U06-GRAM-U06-SOURCE-RELATIVES
% mdparagraph:P248:begin
And set before him the provisions that he had to offer him.
% mdparagraph:P248:end

% mdparagraph:P249:begin
\textbf{Ēstas} is supplied masculine accusative plural, “provisions; food”; \textbf{lęgde} is a supplied third-singular preterite, “placed; set.” \textbf{Him beforan} means “before him,” with the dative person. The following \textbf{ðe} is an object-relative referring to the provisions: “that/\allowbreak{}which he had to offer him.”
% mdparagraph:P249:end

% mdparagraph:P250:begin
Here “had to offer” means \textbf{had available for offering}, not “was obliged to offer.” Tō bēodanne has the learned inflected-infinitive form; bēodan’s “offer” meaning is familiar. The whole availability construction, its relative clause, and its offered-object/\allowbreak{}recipient relation remain supplied. Both occurrences of him refer to Apollonius; the second is not a reflexive “himself.”
% mdparagraph:P250:end

\reserve{5\baselineskip}
% mdparagraph:P251:begin
\textbf{5. Ðā gīet hē wolde, be his mihte, māran ārfæstnesse him gecȳðan;}
% mdparagraph:P251:end

% support: U06-SRC-SUP-KINDNESS
% teaching: NEW-U06-SUPFORM-giet-invariant-source
% teaching: NEW-U06-SUPSENSE-giet-still-source
% teaching: NEW-U06-SUPFORM-be-invariant-source
% teaching: NEW-U06-SUPSENSE-be-according-source
% teaching: NEW-U06-SUPFORM-miht-dat-sg-source
% teaching: NEW-U06-SUPSENSE-miht-power-source
% teaching: NEW-U06-SUPFORM-arfaestnes-acc-sg-source
% teaching: NEW-U06-SUPSENSE-arfaestnes-kindness-source
% teaching: NEW-U06-SUPFORM-gecythan-infinitive-source
% teaching: NEW-U06-SUPSENSE-gecythan-show-source
% mdparagraph:P252:begin
He still wanted, according to his ability, to show him greater kindness.
% mdparagraph:P252:end

% mdparagraph:P253:begin
\textbf{Gīet} means “still; further.” \textbf{Be his mihte}, “according to his ability,” contains dative singular mihte from the noun miht, “ability; power.” It is not a modal verb here. \textbf{Ārfæstnesse} is accusative singular, “kindness”; māran agrees with it as the comparative “greater.” \textbf{Gecȳðan} is a supplied plain infinitive, “show.” In this exact use the thing shown is accusative and its recipient is dative him, Apollonius. The finite wolde expresses the fisherman’s past desire. The prefixed verb and its local government are supplied, not a rule for forming new prefixed verbs.
% mdparagraph:P253:end

\reserve{5\baselineskip}
% mdparagraph:P254:begin
\textbf{6. tōslāt ðā his wǣfels on twā,}
% mdparagraph:P254:end

% support: U06-SRC-SUP-TEARING
% teaching: NEW-U06-SUPFORM-toslitan-pret-3-sg-source
% teaching: NEW-U06-SUPSENSE-toslitan-tear-source
% teaching: NEW-U06-SUPFORM-waefels-acc-sg-source
% teaching: NEW-U06-SUPSENSE-waefels-cloak-source
% teaching: NEW-U06-SUPSENSE-on-divided-into-two-source
% teaching: NEW-U06-SUPSENSE-twegen-divided-in-two-source
% teaching: NEW-U06-GRAM-U06-SOURCE-ELLIPSIS
% mdparagraph:P255:begin
He then tore his cloak in two.
% mdparagraph:P255:end

% mdparagraph:P256:begin
\textbf{Tōslāt} is the supplied third-singular preterite, “tore.” \textbf{Wǣfels} is the accusative singular object, “cloak.” The fisherman remains the subject without a repeated pronoun. Keep \textbf{on twā}, “in two,” as a whole supplied expression. Its usual analysis takes twā as neuter accusative, but that spelling also occurs in feminine forms; it does not identify an unexpressed noun or establish a general rule for motion with on.
% mdparagraph:P256:end

\reserve{5\baselineskip}
% mdparagraph:P257:begin
\textbf{7. and sealde Apollonie ðone healfan dǣl,}
% mdparagraph:P257:end

% support: U06-SRC-SUP-GIFT
% teaching: NEW-U06-SUPFORM-apollonius-dat-sg-source
% teaching: NEW-U06-SUPSENSE-apollonius-person
% teaching: NEW-U06-SUPFORM-healf-acc-m-sg-weak-source
% teaching: NEW-U06-SUPSENSE-healf-half-source
% teaching: NEW-U06-SUPFORM-dael-acc-sg-source
% teaching: NEW-U06-SUPSENSE-dael-part-source
% mdparagraph:P258:begin
And gave Apollonius one half.
% mdparagraph:P258:end

% mdparagraph:P259:begin
\textbf{Apollonie} is the supplied dative singular name, “to Apollonius.” \textbf{Ðone healfan dǣl} means “one half,” literally “the half part.” Dǣl, “part,” is accusative singular; healfan is the weak masculine accusative singular of healf, “half.” This is the familiar giving construction with its recipient and thing given. No second object has to be added.
% mdparagraph:P259:end

\reserve{5\baselineskip}
% mdparagraph:P260:begin
\textbf{8. ðus cweðende:}
% mdparagraph:P260:end

% support: U06-SRC-SUP-SPEECH
% teaching: NEW-U06-SUPFORM-thus-invariant-source
% teaching: NEW-U06-SUPSENSE-thus-thus-source
% teaching: NEW-U06-SUPFORM-cwethan-present-participle-source
% teaching: NEW-U06-GRAM-U06-SOURCE-PARTICIPLE
% mdparagraph:P261:begin
Saying thus.
% mdparagraph:P261:end

% mdparagraph:P262:begin
\textbf{Ðus} means “thus.” \textbf{Cweðende}, “saying,” is a supplied present participle accompanying the fisherman’s action. It introduces his words without serving as a finite past or an imperative. Compare its role with the familiar finite cwæð; no productive participial construction is required here.
% mdparagraph:P262:end

\reserve{5\baselineskip}
% mdparagraph:P263:begin
\textbf{9. “Nim ðæt ic ðē tō sęllanne hæbbe,}
% mdparagraph:P263:end

% support: U06-SRC-SUP-TAKE
% teaching: NEW-U06-SUPSENSE-se-free-relative-what-source
% mdparagraph:P264:begin
Take what I have to give you.
% mdparagraph:P264:end

% mdparagraph:P265:begin
\textbf{Nim} is a singular command. Supply the whole following phrase \textbf{ðæt ic ðē tō sęllanne hæbbe}, “what I have to give you.” Here ðæt introduces a definite free relative, “what/\allowbreak{}that which”; it is not merely the article before an omitted noun in an independently learned rule.
% mdparagraph:P265:end

% mdparagraph:P266:begin
“Have to give” again means \textbf{have available to give}, not obligation. The fisherman is ic; Apollonius is the dative recipient ðē. Sęllanne has the learned inflected-infinitive form. Hæbbe is read here as the familiar first-singular indicative, but its isolated spelling is also shared with subjunctive forms: the letters alone do not uniquely identify its mood. The whole relative and availability construction remains supplied.
% mdparagraph:P266:end

\reserve{5\baselineskip}
% mdparagraph:P267:begin
\textbf{10. and gā intō ðǣre ceastre.”}
% mdparagraph:P267:end

% support: U06-SRC-SUP-GO
% teaching: NEW-U06-SUPFORM-into-invariant-source
% teaching: NEW-U06-SUPSENSE-into-destination-source
% teaching: NEW-U06-SUPFORM-ceaster-dat-sg-source
% teaching: NEW-U06-SUPSENSE-ceaster-city-source
% mdparagraph:P268:begin
And go into the town.
% mdparagraph:P268:end

% mdparagraph:P269:begin
\textbf{Gā} is a singular command. \textbf{Intō ðǣre ceastre} is supplied as “into the town.” Ceastre is dative singular, with matching feminine dative ðǣre. Preserve this actual phrase; the presence of motion does not justify changing it to an accusative or deriving a general rule for every destination. This completes the first spoken direction; do not add the continuation.
% mdparagraph:P269:end

% mdparagraph:P270:begin
\textbf{Source note.} The excerpt is Cook Reader XII, p.167 /\allowbreak{} PDF p.183, editorial lines 3–10. The source has \textbf{læg} without a length mark and \textbf{gīet} with long ī. Line breaks are joined and the attached note cues are moved below. The final closing quotation mark marks the editorial stopping point; it is not printed there in the continuing speech.
% mdparagraph:P270:end

\reserve{7\baselineskip}
% mdparagraph:P271:begin
\subsection*{Cook’s original notes}
% mdparagraph:P271:end

% mdparagraph:P272:begin
The numbered references below are Cook’s, not sections of this unit. The local help above supplies the needed interpretation.
% mdparagraph:P272:end

% source-note: FB167-n5
% mdparagraph:P273:begin
\textbf{P.167, note 5.} fiscere: See 143.
% mdparagraph:P273:end

% source-note: FB167-n6
% mdparagraph:P274:begin
\textbf{P.167, note 6.} ēstas: Lat. epulas.
% mdparagraph:P274:end

% source-note: FB167-n7
% mdparagraph:P275:begin
\textbf{P.167, note 7.} ārfæstnesse: MS. fæstnesse. Lat. pietati.
% mdparagraph:P275:end

\reserve{7\baselineskip}
% mdparagraph:P276:begin
\subsection*{Practice R. The original excerpt}
% mdparagraph:P276:end

% mdparagraph:P277:begin
Use all the printed local help. Identification of learned forms does not require independent recall of the supplied clauses.
% mdparagraph:P277:end

\reserve{5\baselineskip}
\par\begin{minipage}{\linewidth}\raggedright
% mdparagraph:P278:begin
\textbf{R1.}
% mdparagraph:P278:end

% item: NEW-U06-R01
% mdparagraph:P279:begin
Read “The fisherman’s gift and first direction”, using its local support. What prompts the fisherman’s action, what does he divide, and what does he give Apollonius?
% mdparagraph:P279:end

\end{minipage}\par\addvspace{4pt}

\par\begin{minipage}{\linewidth}\raggedright
% mdparagraph:P280:begin
\textbf{R2.}
% mdparagraph:P280:end

% item: NEW-U06-R02
% mdparagraph:P281:begin
Find wolde and its plain infinitive in the source. Parse wolde and explain the fisherman’s intention. Use the supplied meaning of gecȳðan.
% mdparagraph:P281:end

\end{minipage}\par\addvspace{4pt}

\par\begin{minipage}{\linewidth}\raggedright
% mdparagraph:P282:begin
\textbf{R3.}
% mdparagraph:P282:end

% item: NEW-U06-R03
% mdparagraph:P283:begin
Identify tō bēodanne and tō sęllanne. Give the familiar plain infinitive behind each. Explain the source’s “have to offer/\allowbreak{}give” phrases using the printed support; do not translate them as modern English “must”.
% mdparagraph:P283:end

\end{minipage}\par\addvspace{4pt}

\par\begin{minipage}{\linewidth}\raggedright
% mdparagraph:P284:begin
\textbf{R4.}
% mdparagraph:P284:end

% item: NEW-U06-R04
% mdparagraph:P285:begin
Find the source’s two singular commands. Parse each and give its plural command form from the taught system.
% mdparagraph:P285:end

\end{minipage}\par\addvspace{4pt}

\par\begin{minipage}{\linewidth}\raggedright
% mdparagraph:P286:begin
\textbf{R5.}
% mdparagraph:P286:end

% item: NEW-U06-R05
% mdparagraph:P287:begin
Use the notes to distinguish the participial cweðende from the finite past cwæð and the command gā. Which form introduces the fisherman’s actual words?
% mdparagraph:P287:end

\end{minipage}\par\addvspace{4pt}

\par\begin{minipage}{\linewidth}\raggedright
% mdparagraph:P288:begin
\textbf{R6.}
% mdparagraph:P288:end

% item: NEW-U06-R06
% mdparagraph:P289:begin
What does the fisherman say he has, and what first direction does he give? Explain the force of the phrase translated “what I have” using the local support.
% mdparagraph:P289:end

\end{minipage}\par\addvspace{4pt}

\par\begin{minipage}{\linewidth}\raggedright
% mdparagraph:P290:begin
\textbf{R7.}
% mdparagraph:P290:end

% item: NEW-U06-R07
% mdparagraph:P291:begin
Identify the source’s tō phrases that contain an inflected infinitive and distinguish them from tō his hūse and intō ðǣre ceastre. State the locally supplied meanings of the destination phrases without deriving a general motion-preposition rule.
% mdparagraph:P291:end

\end{minipage}\par\addvspace{4pt}

\par\begin{minipage}{\linewidth}\raggedright
% mdparagraph:P292:begin
\textbf{R8.}
% mdparagraph:P292:end

% item: NEW-U06-R08
% mdparagraph:P293:begin
Translate the whole 76-word source excerpt with its supports. Stop after the fisherman’s first direction; do not add the continuation of his speech.
% mdparagraph:P293:end

\end{minipage}\par\addvspace{4pt}

\reserve{7\baselineskip}
% mdparagraph:P294:begin
\section*{8. Cumulative practice}
% mdparagraph:P294:end

\reserve{7\baselineskip}
% mdparagraph:P295:begin
\subsection*{Practice X. Combining the constructions}
% mdparagraph:P295:end

% mdparagraph:P296:begin
The final task admits more than one coherent answer. Keep to its lexical bank and the constructions taught here; the four sentences should form a connected account.
% mdparagraph:P296:end

\reserve{5\baselineskip}
\par\begin{minipage}{\linewidth}\raggedright
% mdparagraph:P297:begin
\textbf{X1.}
% mdparagraph:P297:end

% item: NEW-U06-X01
% mdparagraph:P298:begin
Translate “We two wanted to give our food to the birds. We two gave it to them.” Retain the two-person possessor and the grammatical gender of “food”.
% mdparagraph:P298:end

\end{minipage}\par\addvspace{4pt}

\par\begin{minipage}{\linewidth}\raggedright
% mdparagraph:P299:begin
\textbf{X2.}
% mdparagraph:P299:end

% item: NEW-U06-X02
% mdparagraph:P300:begin
Translate “The old man cannot carry the great ship. Help him.” Address several companions in the command.
% mdparagraph:P300:end

\end{minipage}\par\addvspace{4pt}

\par\begin{minipage}{\linewidth}\raggedright
% mdparagraph:P301:begin
\textbf{X3.}
% mdparagraph:P301:end

% item: NEW-U06-X03
% mdparagraph:P302:begin
Translate “He came to help us. We wanted to give him food. He sat and ate.” Use the taught elementary purpose and volitional constructions.
% mdparagraph:P302:end

\end{minipage}\par\addvspace{4pt}

\par\begin{minipage}{\linewidth}\raggedright
% mdparagraph:P303:begin
\textbf{X4.}
% mdparagraph:P303:end

% item: NEW-U06-X04
% mdparagraph:P304:begin
Compare the statement Gē gāð with the command Gāð. Explain the shared spelling and the difference in use. Then compare Hē dēð ðæt weorc with Dō ðæt weorc.
% mdparagraph:P304:end

\end{minipage}\par\addvspace{4pt}

\par\begin{minipage}{\linewidth}\raggedright
% mdparagraph:P305:begin
\textbf{X5.}
% mdparagraph:P305:end

% item: NEW-U06-X05
% mdparagraph:P306:begin
Write four connected Old English sentences about a man and two companions: one past arrival with an explicitly stated purpose, one statement of willingness, one direct command addressed to both companions, and one completed past action. Use cuman, helpan, willan and beran, with mann, gefēra and męte and any needed taught articles, numerals or pronouns. Label the three new constructions.
% mdparagraph:P306:end

\end{minipage}\par\addvspace{4pt}

\clearpage
% mdparagraph:P307:begin
\section*{9. Sources and form choices}
% mdparagraph:P307:end

\begingroup\fontsize{10}{12.8}\selectfont
\setlength{\parskip}{1.5pt}
% mdparagraph:P308:begin
The core forms follow Albert S. Cook, \textbf{A First Book in Old English: Grammar, Reader, Notes, and Vocabulary}, second edition, revised and enlarged (Ginn, 1895). The compact tables explicitly identify the three persons covered by each shared plural row and apply the given endings to the selected stems. Such paradigm completion is not a claim that each cell is separately quoted in continuous prose.
% mdparagraph:P308:end

% mdparagraph:P309:begin
Joseph Wright and Elizabeth Mary Wright, \textbf{Old English Grammar}, second edition (Oxford, 1914), corroborates the forms and variants; the precise witan and magan variant locators follow below. R. D. Fulk, \textbf{An Introductory Grammar of Old English with an Anthology of Readings}, author-corrected open-access reprint (2020), supplies the cited distinctions concerning defects, j-presents, the bēon series, and infinitival use. \textbf{Bosworth–Toller, An Anglo-Saxon Dictionary}, main dictionary (1898), supplies the individually cited lexical senses and constructions; page letters a and b identify columns.
% mdparagraph:P309:end

% mdparagraph:P310:begin
All PDF page numbers in this unit are one-based. In the selected Cook scan the PDF page is printed page +16; in Wright, +20; in Fulk, +12. Dictionary facsimiles are cited by printed page and column rather than by an invented PDF position.
% mdparagraph:P310:end

% mdparagraph:P311:begin
The examples and “Food and water for the journey” are new compositions. The original fisherman excerpt preserves its source wording, with only the declared closing quotation mark added. Its supplied nouns, participle, prepositions, and clause constructions remain local help. Earlier participles and other reference-only forms do not become command or infinitive systems through resemblance in spelling.
% mdparagraph:P311:end

% mdparagraph:P312:begin
\textbf{Local reading help.} Compare Cook §§87a and 87d, pp.51–52 /\allowbreak{} PDF pp.67–68, for relatives; §164b, p.94 /\allowbreak{} PDF p.110, for offering and giving; and §§194–200, pp.104–106 /\allowbreak{} PDF pp.120–122, for the verbal forms and uses. The availability and relative phrases in the reading remain supplied as a whole.
% mdparagraph:P312:end

% mdparagraph:P313:begin
The supplied lexical help follows the corresponding entries in Cook’s \textbf{Vocabulary}, whose headword spellings are retained in these references: æt (p.274 /\allowbreak{} PDF p.290); ār-fæstnes (p.276 /\allowbreak{} PDF p.292); be (p.277 /\allowbreak{} PDF p.293); be-foran; bēodan (p.278 /\allowbreak{} PDF p.294); ceaster (p.281 /\allowbreak{} PDF p.297); dǣl; cuman (p.282 /\allowbreak{} PDF p.298); ēst (p.285 /\allowbreak{} PDF p.301); fiscere (p.287 /\allowbreak{} PDF p.303); fōt (p.289 /\allowbreak{} PDF p.305); ge-cȳðan (p.291 /\allowbreak{} PDF p.307); geong (p.293 /\allowbreak{} PDF p.309); gīet (p.296 /\allowbreak{} PDF p.312); healf (p.298 /\allowbreak{} PDF p.314); hūs (p.300 /\allowbreak{} PDF p.316); in-tō (p.301 /\allowbreak{} PDF p.317); lęcgean (p.302 /\allowbreak{} PDF p.318); mid; miht; mild-heortnes (p.305 /\allowbreak{} PDF p.321); sōna (p.313 /\allowbreak{} PDF p.329); tō; tō-slītan (p.316 /\allowbreak{} PDF p.332); ðus (p.319 /\allowbreak{} PDF p.335); ūp-ā-hębban; wǣfels (p.320 /\allowbreak{} PDF p.336).
% mdparagraph:P313:end

% mdparagraph:P314:begin
\textbf{Witan variants.} Wright, §540, printed p.283 /\allowbreak{} PDF p.303.
% mdparagraph:P314:end

% mdparagraph:P315:begin
\textbf{Magan variants.} Wright, §§544–545, printed p.285 /\allowbreak{} PDF p.305.
% mdparagraph:P315:end

\endgroup
\end{document}
